\subsection{Per-result source and status concordance}\label{app:source-status}
The complete statement text is in the referenced result and, verbatim as TeX, in \path{S6_07_PAPER_source_table.csv}. The table below gives one row per printed statement or explicitly recorded setup. T, K and A are the binding statuses; \emph{text} denotes definitions or setup. The credit checks and precise wording changes are retained per row in that CSV, with the full source-to-print diffs in the accompanying records. Source filenames are relative to the supplied bundles.
\begingroup\small\setlength{\tabcolsep}{3pt}
\begin{longtable}{@{}L{.18\linewidth}L{.48\linewidth}L{.29\linewidth}@{}}
\toprule Result / status & Statement locator and source & Kind / proof\\\midrule\endhead
FAM \nolinkurl{conventions}\newline text & Statement~\ref{def:conventions}. Base, lattice, periods, clockwise deck generators and transported reference.\newline \path{returns/S6_02_FAMILY/sections/conventions.tex} & Definition\newline §1\\
FAM \nolinkurl{A.1.1}\newline K & Statement~\ref{fam:A.1.1}. Intrinsic pencil, canonical ring and base preservation.\newline \path{returns/S6_02_FAMILY/sections/AB.tex} & New proof\newline Appendix B\\
DGKXv1 Theorem 7.2; thm:\allowbreak{}all-maps-factor\newline K & Statement~\ref{dgkx:all-maps-factor}. Every holomorphic map to P¹ factors uniquely through f.\newline \path{papers/arXiv_2610.00250v1_Du_Guo_Kusakabe_Xie.tex} & New proof\newline Appendix B (source re-derivation)\\
FAM \nolinkurl{B.1.1}\newline T & Statement~\ref{fam:B.1.1}. Elliptic-surface quotient and Mordell--Weil generator.\newline \path{returns/S6_02_FAMILY/sections/AB.tex} & Theorem\newline §1\\
FAM \nolinkurl{B.1.2}\newline T & Statement~\ref{fam:B.1.2}. The exact admissible domain for the reference quotient.\newline \path{returns/S6_02_FAMILY/sections/AB.tex} & Theorem\newline §1\\
FAM \nolinkurl{B.1.3}\newline T & Statement~\ref{fam:B.1.3}. The distinguished elliptic logarithm is a zero of the Weierstrass function.\newline \path{returns/S6_02_FAMILY/sections/AB.tex} & New proof\newline §1\\
FAM \nolinkurl{B.1.4}\newline T & Statement~\ref{fam:B.1.4}. The transported reference and the square relation between parameters.\newline \path{returns/S6_02_FAMILY/sections/AB.tex} & Theorem\newline §1\\
FAM \nolinkurl{B.2.3}\newline T & Statement~\ref{fam:B.2.3}. Correction of the global descent claim in Engel Remark 8.1.\newline \path{returns/S6_02_FAMILY/sections/AB.tex} & Correction\newline §1\\
FAM \nolinkurl{A.5.1}\newline K & Statement~\ref{fam:A.5.1}. The known exact identification for the standard data.\newline \path{returns/S6_02_FAMILY/sections/AB.tex} & New proof\newline Appendices A and B\\
FAM \nolinkurl{B.2.1}\newline T & Statement~\ref{fam:B.2.1}. Complete invariants for all admissible discrete data.\newline \path{returns/S6_02_FAMILY/sections/AB.tex} & Theorem\newline §2\\
FAM \nolinkurl{B.2.2}\newline T & Statement~\ref{fam:B.2.2}. Two classes at a fixed Engel parameter.\newline \path{returns/S6_02_FAMILY/sections/AB.tex} & Theorem\newline §2\\
FAM \nolinkurl{B.3.1}\newline T & Statement~\ref{fam:B.3.1}. Identification subgroup for the stated higher-multiplicity model.\newline \path{returns/S6_02_FAMILY/sections/AB.tex} & Theorem\newline §2\\
PUSH \nolinkurl{E.1.1}\newline A & Statement~\ref{push:E.1.1}. The coordinate and primitive-circle dictionary.\newline \path{returns/S6_02_PUSH/sections/E_push.tex} & Theorem\newline Repository:\allowbreak{} S6\_\allowbreak{}02\_\allowbreak{}PUSH E.1.1\\
PUSH \nolinkurl{E.1.2}\newline A & Statement~\ref{push:E.1.2}. The higher-multiplicity parameter dictionary.\newline \path{returns/S6_02_PUSH/sections/E_push.tex} & Theorem\newline Repository:\allowbreak{} S6\_\allowbreak{}02\_\allowbreak{}PUSH E.1.2\\
FAM \nolinkurl{A.1.2}\newline T & Statement~\ref{fam:A.1.2}. Normalized finite base changes and product covers.\newline \path{returns/S6_02_FAMILY/sections/AB.tex} & Theorem\newline Appendix B\\
FAM \nolinkurl{A.2.1}\newline K & Statement~\ref{fam:A.2.1}. Centralizer, finite norms and the parity obstruction.\newline \path{returns/S6_02_FAMILY/sections/AB.tex} & New proof\newline Appendix B\\
FAM \nolinkurl{A.3.1}\newline K & Statement~\ref{fam:A.3.1}. Coherent even shifts and the proper quotient family.\newline \path{returns/S6_02_FAMILY/sections/AB.tex} & New proof\newline Appendix B\\
FAM \nolinkurl{A.4.1}\newline T & Statement~\ref{fam:A.4.1}. Albanese lattices and the generic two-invariant comparison.\newline \path{returns/S6_02_FAMILY/sections/AB.tex} & Theorem\newline Appendix B\\
FAM \nolinkurl{A.4.2}\newline K & Statement~\ref{fam:A.4.2}. Nonvanishing of the Kodaira--Spencer class.\newline \path{returns/S6_02_FAMILY/sections/AB.tex} & New proof\newline Appendix B\\
WORKBENCH \nolinkurl{pshift:lem:centralizer}\newline K & Statement~\ref{pshift:lem:centralizer}. Complete rational centralizer and middle-family linear parts.\newline \path{workbench/period_shift_classification.tex} & New proof\newline Appendix A\\
WORKBENCH \nolinkurl{pshift:thm:classification}\newline K & Statement~\ref{pshift:thm:classification}. Exact global identification of the retained period parameter.\newline \path{workbench/period_shift_classification.tex} & New proof\newline Appendix A\\
FAM \nolinkurl{C.1.1}\newline T & Statement~\ref{fam:C.1.1}. Theta formula, distinguished Weierstrass zero and special values.\newline \path{returns/S6_02_FAMILY/sections/CD.tex} & Theorem\newline §3\\
FAM \nolinkurl{C.1.2}\newline T & Statement~\ref{fam:C.1.2}. Special-value evaluations and the period formula.\newline \path{returns/S6_02_FAMILY/sections/CD.tex} & Computation\newline §3\\
ARI \nolinkurl{C.1.1}\newline T & Statement~\ref{ari:C.1.1}. distribution under isogenies\newline \path{returns/S6_05_ARITHMETIC/sections/C_special.tex} & Theorem\newline §3\\
ARI \nolinkurl{C.1.2}\newline T & Statement~\ref{ari:C.1.2}. the phases and the two periods\newline \path{returns/S6_05_ARITHMETIC/sections/C_special.tex} & Theorem\newline §3\\
ARI \nolinkurl{C.2.1}\newline T & Statement~\ref{ari:C.2.1}. exact nomes and continuation radius\newline \path{returns/S6_05_ARITHMETIC/sections/C_special.tex} & Theorem\newline §3\\
ARI \nolinkurl{C.4.2}\newline T & Statement~\ref{ari:C.4.2}. the full division identity\newline \path{returns/S6_05_ARITHMETIC/sections/C_special.tex} & Theorem\newline §3\\
ARI \nolinkurl{C.4.3}\newline T & Statement~\ref{ari:C.4.3}. algebraicity and explicit fields for every order\newline \path{returns/S6_05_ARITHMETIC/sections/C_special.tex} & Theorem\newline §3\\
ARI \nolinkurl{C.4.4}\newline T & Statement~\ref{ari:C.4.4}. the five values and all orders through eight\newline \path{returns/S6_05_ARITHMETIC/sections/C_special.tex} & Theorem\newline §3\\
ARI \nolinkurl{C.8.6}\newline T & Statement~\ref{ari:C.8.6}. the Mahler-measure identities\newline \path{returns/S6_05_ARITHMETIC/sections/C_integral.tex} & Theorem\newline §3\\
ARI \nolinkurl{C.8.7}\newline T & Statement~\ref{ari:C.8.7}. the two complex beta integrals\newline \path{returns/S6_05_ARITHMETIC/sections/C_integral.tex} & Theorem\newline §3\\
ARI \nolinkurl{C.8.8}\newline A & Statement~\ref{ari:C.8.8}. the signed Gaussian identity\newline \path{returns/S6_05_ARITHMETIC/sections/C_integral.tex} & Theorem\newline §3 (mechanism only; repository proof)\\
ARI \nolinkurl{C.8.9}\newline A & Statement~\ref{ari:C.8.9}. C.8.9; Corrected sign and measure normalization\newline \path{returns/S6_05_ARITHMETIC/sections/C_measure.tex} & Theorem\newline §3 (mechanism only; repository proof)\\
ARI \nolinkurl{C.8.10}\newline A & Statement~\ref{ari:C.8.10}. the correctly normalized equidistribution\newline \path{returns/S6_05_ARITHMETIC/sections/C_measure.tex} & Theorem\newline §3 (mechanism only; repository proof)\\
ARI \nolinkurl{C.8.11}\newline A & Statement~\ref{ari:C.8.11}. the two first-derivative identities\newline \path{returns/S6_05_ARITHMETIC/sections/C_measure.tex} & Theorem\newline §3 (mechanism only; repository proof)\\
FAM \nolinkurl{C.2.3}\newline T & Statement~\ref{fam:C.2.3}. The exact maximum and the radius of the parameter disc.\newline \path{returns/S6_02_FAMILY/sections/CD.tex} & Computation\newline §4\\
PUSH \nolinkurl{E.4.1}\newline T & Statement~\ref{push:E.4.1}. independent certification of every box\newline \path{returns/S6_02_PUSH/sections/E4_push.tex} & Computation\newline §4\\
ARI \nolinkurl{C.3.1}\newline T & Statement~\ref{ari:C.3.1}. fourth order at the maximum\newline \path{returns/S6_05_ARITHMETIC/sections/C_special.tex} & Theorem\newline §4\\
ARI \nolinkurl{C.8.1}\newline T & Statement~\ref{ari:C.8.1}. three uniform bounds on the elliptic norm\newline \path{returns/S6_05_ARITHMETIC/sections/C_analytic.tex} & Theorem\newline §4\\
ARI \nolinkurl{C.8.2}\newline T & Statement~\ref{ari:C.8.2}. an explicit local maximum on a larger disc\newline \path{returns/S6_05_ARITHMETIC/sections/C_analytic.tex} & Theorem\newline §4\\
ARI \nolinkurl{C.8.3}\newline T & Statement~\ref{ari:C.8.3}. the comparison of the two local discs\newline \path{returns/S6_05_ARITHMETIC/sections/C_analytic.tex} & Theorem\newline §4\\
ARI \nolinkurl{C.8.4}\newline T & Statement~\ref{ari:C.8.4}. the first three normalized duplication steps\newline \path{returns/S6_05_ARITHMETIC/sections/C_analytic.tex} & Theorem\newline §4\\
ARI \nolinkurl{C.8.5}\newline T & Statement~\ref{ari:C.8.5}. Maximum by explicit analytic estimates\newline \path{returns/S6_05_ARITHMETIC/sections/C_analytic.tex} & Theorem\newline §4\\
ARI \nolinkurl{C.8.15}\newline A & Statement~\ref{ari:C.8.15}. every critical point and every level set\newline \path{returns/S6_05_ARITHMETIC/sections/C_critical.tex} & Theorem\newline §4 (mechanism only; repository proof)\\
FAM C.1 setup\newline text & \nolinkurl{periods:theta-definitions}. C.1 setup\newline \path{returns/S6_02_FAMILY/sections/CD.tex} & Definition/\allowbreak{}record\newline Setup\\
FAM C.1.2 artifact-availability record\newline text & \nolinkurl{periods:special-computation-availability}. C.1.2 artifact-availability record\newline \path{returns/S6_02_FAMILY/sections/CD.tex} & Artifact record\newline Setup\\
ARI C.1 setup; C.8 critical setup\newline text & \nolinkurl{periods:norm-definition}. C.1 setup; C.8 critical setup\newline \path{returns/S6_05_ARITHMETIC/sections/C_special.tex}\newline \path{returns/S6_05_ARITHMETIC/sections/C_critical.tex} & Definition/\allowbreak{}record\newline Setup\\
FAM C.2.1 proof, eq:\allowbreak{}cd-norm\newline text & \nolinkurl{fam:eq:cd-norm}. C.2.1 proof, eq:\allowbreak{}cd-norm\newline \path{returns/S6_02_FAMILY/sections/CD.tex} & Definition/\allowbreak{}record\newline Setup\\
ARI C.4 setup and local-height notation in C.4.1\newline text & \nolinkurl{periods:division-model}. C.4 setup and local-height notation in C.4.1\newline \path{returns/S6_05_ARITHMETIC/sections/C_special.tex} & Definition/\allowbreak{}record\newline Setup\\
ARI C.8 measure conventions\newline text & \nolinkurl{periods:measure-conventions}. C.8 measure conventions\newline \path{returns/S6_05_ARITHMETIC/sections/C_measure.tex} & Definition/\allowbreak{}record\newline Setup\\
FAM;PUSH FAM C.2.3; PUSH E.4.1--3; REVIEW\_\allowbreak{}S6\_\allowbreak{}02\_\allowbreak{}FAMILY §2\newline text & \nolinkurl{periods:certificate-record}. FAM C.2.3; PUSH E.4.1--3; REVIEW\_\allowbreak{}S6\_\allowbreak{}02\_\allowbreak{}FAMILY §2\newline \path{returns/S6_02_FAMILY/sections/CD.tex}\newline \path{returns/S6_02_PUSH/sections/E4_push.tex}\newline \path{status/REVIEW_S6_02_FAMILY.md} & Computation record\newline Setup\\
ARI C.3 setup\newline text & \nolinkurl{periods:local-expansion-constants}. C.3 setup\newline \path{returns/S6_05_ARITHMETIC/sections/C_special.tex} & Definition/\allowbreak{}record\newline Setup\\
ARI C.8 analytic setup\newline text & \nolinkurl{periods:ladder-modular-conventions}. C.8 analytic setup\newline \path{returns/S6_05_ARITHMETIC/sections/C_analytic.tex} & Definition/\allowbreak{}record\newline Setup\\
INV \nolinkurl{notation-cotangent}\newline text & Statement~\ref{inv:notation-cotangent}. Cotangent notation\newline \path{returns/S6_01_INVARIANTS/fragments/part_A.tex} & Definition\newline definition\\
INV \nolinkurl{A1.0}\newline T & Statement~\ref{inv:A1.0}. The structure-sheaf input in the precise form used below\newline \path{returns/S6_01_INVARIANTS/fragments/part_A.tex} & New proof\newline main text\\
INV \nolinkurl{A1.1}\newline T & Statement~\ref{inv:A1.1}. Local freeness and the determinant line\newline \path{returns/S6_01_INVARIANTS/fragments/part_A.tex} & Theorem\newline main text\\
INV \nolinkurl{A1.2}\newline T & Statement~\ref{inv:A1.2}. The logarithmic direct images by vertical twisting and duality\newline \path{returns/S6_01_INVARIANTS/fragments/part_A.tex} & Theorem\newline main text\\
INV \nolinkurl{notation-period-branch}\newline text & Statement~\ref{inv:notation-period-branch}. The order-four branch\newline \path{returns/S6_01_INVARIANTS/fragments/part_A.tex} & Definition\newline definition\\
INV \nolinkurl{A2.1}\newline T & Statement~\ref{inv:A2.1}. The complete raw period determinant\newline \path{returns/S6_01_INVARIANTS/fragments/part_A.tex} & Theorem\newline main text\\
INV \nolinkurl{A2.2}\newline T & Statement~\ref{inv:A2.2}. The evaluated scalar and the two local leading terms\newline \path{returns/S6_01_INVARIANTS/fragments/part_A.tex} & Theorem\newline main text\\
INV \nolinkurl{A2.4}\newline T & Statement~\ref{inv:A2.4}. The global rank and Lead 2\newline \path{returns/S6_01_INVARIANTS/fragments/part_A.tex} & Theorem\newline main text\\
INV \nolinkurl{A1.3}\newline T & Statement~\ref{inv:A1.3}. All absolute cotangent direct images\newline \path{returns/S6_01_INVARIANTS/fragments/part_A.tex} & Theorem\newline main text\\
INV \nolinkurl{A3.1}\newline T & Statement~\ref{inv:A3.1}. The complete Hodge table for every admissible parameter\newline \path{returns/S6_01_INVARIANTS/fragments/part_A.tex} & Theorem\newline main text\\
INV \nolinkurl{notation-hodge-input}\newline text & Statement~\ref{inv:notation-hodge-input}. The formal Hodge parameter\newline \path{returns/S6_01_INVARIANTS/fragments/part_C.tex} & Definition\newline definition\\
INV \nolinkurl{C1.1}\newline T & Statement~\ref{inv:C1.1}. The complete Frölicher differential in terms of h.\newline \path{returns/S6_01_INVARIANTS/fragments/part_C.tex} & Theorem\newline main text\\
INV \nolinkurl{C2.1}\newline T & Statement~\ref{inv:C2.1}. Bott--Chern and Aeppli tables as functions of h.\newline \path{returns/S6_01_INVARIANTS/fragments/part_C.tex} & Theorem\newline main text\\
INV \nolinkurl{C3.3}\newline A & Statement~\ref{inv:C3.3}. The complete Bott--Chern algebra, for every h.\newline \path{returns/S6_01_INVARIANTS/fragments/part_C.tex} & Theorem\newline repository:\allowbreak{} \path{returns/S6_01_INVARIANTS/fragments/part_C.tex}\\
INV \nolinkurl{B1.1}\newline T & Statement~\ref{inv:B1.1}. The relative cohomology comparison\newline \path{returns/S6_01_INVARIANTS/fragments/part_B.tex} & Theorem\newline main text\\
INV \nolinkurl{B2.1}\newline T & Statement~\ref{inv:B2.1-canonical}. Canonical bundles for all coprime multiplicities\newline \path{returns/S6_01_INVARIANTS/fragments/part_B.tex} & Theorem\newline main text\\
INV \nolinkurl{B1.2}\newline A & Statement~\ref{inv:B1.2}. All admissible twists have the same Hodge numbers\newline \path{returns/S6_01_INVARIANTS/fragments/part_B.tex} & Theorem\newline repository:\allowbreak{} \path{returns/S6_01_INVARIANTS/fragments/part_B.tex}\\
INV \nolinkurl{B2.1}\newline A & Statement~\ref{inv:B2.1-hodge}. Hodge numbers for all coprime multiplicities\newline \path{returns/S6_01_INVARIANTS/fragments/part_B.tex} & Theorem\newline repository:\allowbreak{} \path{returns/S6_01_INVARIANTS/fragments/part_B.tex}\\
INV \nolinkurl{B3.1}\newline T & Statement~\ref{inv:B3.1}. Connected cyclic covers and the cusp comparison\newline \path{returns/S6_01_INVARIANTS/fragments/part_B.tex} & Theorem\newline main text\\
INV \nolinkurl{B3.2}\newline T & Statement~\ref{inv:B3.2}. Cotangent direct images and Hodge numbers of the covers\newline \path{returns/S6_01_INVARIANTS/fragments/part_B.tex} & Theorem\newline main text\\
INV \nolinkurl{B3.3}\newline T & Statement~\ref{inv:B3.3}. Canonical rings and rational topology of the covers\newline \path{returns/S6_01_INVARIANTS/fragments/part_B.tex} & Theorem\newline main text\\
INV \nolinkurl{C3.5}\newline A & Statement~\ref{inv:C3.5}. Metric obstructions that do not require h.\newline \path{returns/S6_01_INVARIANTS/fragments/part_C.tex} & Theorem\newline repository:\allowbreak{} \path{returns/S6_01_INVARIANTS/fragments/part_C.tex}\\
INV \nolinkurl{C3.7}\newline A & Statement~\ref{inv:C3.7}. The fibre obstruction to SKT and astheno-Kähler metrics.\newline \path{returns/S6_01_INVARIANTS/fragments/part_C.tex} & Theorem\newline repository:\allowbreak{} \path{returns/S6_01_INVARIANTS/fragments/part_C.tex}\\
INV \nolinkurl{A4.1}\newline T & Statement~\ref{inv:A4.1}. All local modules needed to distinguish degrees one and two\newline \path{returns/S6_01_INVARIANTS/fragments/part_A.tex} & Theorem\newline main text\\
INV \nolinkurl{A4.2}\newline T & Statement~\ref{inv:A4.2}. The claimed degree-one torsion along the cusp\newline \path{returns/S6_01_INVARIANTS/fragments/part_A.tex} & Correction\newline main text\\
FAM \nolinkurl{A.5.2}\newline K & Statement~\ref{fam:A.5.2}. Tangent cohomology and the Kuranishi line\newline \path{returns/S6_02_FAMILY/sections/AB.tex} & New proof\newline Appendix B, tangent cohomology subsection\\
CDP \nolinkurl{notation-integral-twists}\newline text & Statement~\ref{cdp:notation-integral-twists}. The integral Picard normalization\newline \path{returns/S6_03_CDP/sections/part_C5.tex} & Definition\newline definition\\
CDP \nolinkurl{C5b}\newline T & Statement~\ref{cdp:C5b}. Integral twists and the normalization-map exception\newline \path{returns/S6_03_CDP/sections/part_C5.tex} & Theorem\newline main text\\
PUSH \nolinkurl{E.2.1}\newline T & Statement~\ref{push:E.2.1}. Tangent cohomology and local miniversality for coprime multiplicities\newline \path{returns/S6_02_PUSH/sections/E_push.tex} & Theorem\newline main text\\
FAM F4 consequence of A.5.2; binding row 6d\newline T & Statement~\ref{fam:A.5.2-corollary}. Algebraic dimension near the standard family\newline \path{status/REVIEW_S6_02_FAMILY.md} & Theorem\newline main text\\
FAM \nolinkurl{B.4.1}\newline K & Statement~\ref{fam:B.4.1}. The full automorphism group is C*.\newline \path{returns/S6_02_FAMILY/sections/AB.tex} & New proof\newline Appendix B, re-derivation of FAM B.4.1\\
FAM \nolinkurl{B.4.2}\newline T & Statement~\ref{fam:B.4.2}. Conjugation sends c to 1−conjugate(c), and the real members have split extended automorphism group.\newline \path{returns/S6_02_FAMILY/sections/AB.tex} & Theorem\newline §7\\
ARI \nolinkurl{C.2.2}\newline T & Statement~\ref{ari:C.2.2}. Real structures occur exactly at nonzero real Albanese nome q2.\newline \path{returns/S6_05_ARITHMETIC/sections/C_special.tex} & Theorem\newline §7\\
MAPS A.3 (opening conventions)\newline text & Statement~\ref{maps:A.3.setup}. Cusp double curves and the real circle action.\newline \path{returns/S6_04_MAPS_MANIFOLDS/sections/real_forms.tex}\newline \path{returns/S6_02_FAMILY/sections/AB.tex} & Definition\newline §7\\
MAPS \nolinkurl{A.3.1}\newline A & Statement~\ref{maps:A.3.1}. There are two holomorphic-conjugacy classes of antiholomorphic involutions.\newline \path{returns/S6_04_MAPS_MANIFOLDS/sections/real_forms.tex} & Theorem\newline §7; proof in repository:\allowbreak{} MAPS A.3.1\\
MAPS A.3, equations (A.3.1)--(A.3.2)\newline text & Statement~\ref{maps:A.3.representatives}. Explicit roots, branches and translating vectors for the two real representatives.\newline \path{returns/S6_04_MAPS_MANIFOLDS/sections/real_forms.tex} & Definition\newline §7\\
MAPS \nolinkurl{A.3.2}\newline A & Statement~\ref{maps:A.3.2}. Exact real fixed sets on the ordinary, multiple and cusp fibres.\newline \path{returns/S6_04_MAPS_MANIFOLDS/sections/real_forms.tex} & Theorem\newline §7; proof in repository:\allowbreak{} MAPS A.3.2\\
MAPS \nolinkurl{A.3.3}\newline A & Statement~\ref{maps:A.3.3}. Every nonempty real fixed manifold is S3, with the specified circle action.\newline \path{returns/S6_04_MAPS_MANIFOLDS/sections/real_forms.tex} & Theorem\newline §7; proof in repository:\allowbreak{} MAPS A.3.3\\
FAM F (opening conventions)\newline text & Statement~\ref{fam:F.setup}. The indefinite period domain and the convention for proper discontinuity.\newline \path{returns/S6_02_FAMILY/sections/F.tex} & Definition\newline §8\\
FAM \nolinkurl{F.1.1}\newline T & Statement~\ref{fam:F.1.1}. The monodromy group is an explicit index-12 character kernel.\newline \path{returns/S6_02_FAMILY/sections/F.tex} & Theorem\newline §8\\
FAM \nolinkurl{F.1.2}\newline T & Statement~\ref{fam:F.1.2}. The centre, extension cocycle and Zariski closures are explicit.\newline \path{returns/S6_02_FAMILY/sections/F.tex} & Theorem\newline §8\\
FAM F.2 (definition preceding F.2.1)\newline text & Statement~\ref{fam:F.2.setup}. The invariant function q\_\allowbreak{}d and its two nonzero sign regions.\newline \path{returns/S6_02_FAMILY/sections/F.tex} & Definition\newline §8\\
FAM \nolinkurl{F.2.1}\newline T & Statement~\ref{fam:F.2.1}. The action is not properly discontinuous on the entire domain.\newline \path{returns/S6_02_FAMILY/sections/F.tex} & Counterexample\newline §8\\
FAM \nolinkurl{F.2.2}\newline T & Statement~\ref{fam:F.2.2}. The complement of q\_\allowbreak{}d=0 is the largest invariant proper-discontinuity domain.\newline \path{returns/S6_02_FAMILY/sections/F.tex} & Theorem\newline §8\\
FAM \nolinkurl{F.2.3}\newline T & Statement~\ref{fam:F.2.3}. All admissible period images lie in a global chart of the positive region.\newline \path{returns/S6_02_FAMILY/sections/F.tex} & Theorem\newline §8\\
FAM \nolinkurl{F.3.1}\newline T & Statement~\ref{fam:F.3.1}. The period curve is an injective immersion and has a proper embedded extension.\newline \path{returns/S6_02_FAMILY/sections/F.tex} & Theorem\newline §8\\
FAM \nolinkurl{F.3.2}\newline T & Statement~\ref{fam:F.3.2}. The finite limits lie in the domain, while the cusp limit lies outside it.\newline \path{returns/S6_02_FAMILY/sections/F.tex} & Theorem\newline §8\\
FAM \nolinkurl{F.3.3}\newline T & Statement~\ref{fam:F.3.3}. The determinant of the derivative is generically nonzero and vanishes at z1.\newline \path{returns/S6_02_FAMILY/sections/F.tex} & Theorem\newline §8\\
FAM \nolinkurl{D.1.1}\newline T & Statement~\ref{fam:D.1.1}. The first projection descends exactly at the six-division points.\newline \path{returns/S6_02_FAMILY/sections/CD.tex} & Theorem\newline §9\\
FAM \nolinkurl{D.1.2}\newline T & Statement~\ref{fam:D.1.2}. The torsion locus is countable and dense, with exact intersection degree.\newline \path{returns/S6_02_FAMILY/sections/CD.tex} & Theorem\newline §9\\
FAM \nolinkurl{D.1.3}\newline T & Statement~\ref{fam:D.1.3}. An unrestricted elliptic-quotient wording is refuted at t=32/\allowbreak{}5.\newline \path{returns/S6_02_FAMILY/sections/CD.tex} & Counterexample\newline §9\\
FAM D.1.6 (proof, Hermitian-form convention)\newline text & Statement~\ref{fam:D.1.Hermitian}. The transferred alternating form and its explicit Hermitian convention.\newline \path{returns/S6_02_FAMILY/sections/CD.tex} & Definition\newline §9\\
FAM \nolinkurl{D.1.4}\newline T & Statement~\ref{fam:D.1.4}. The six-coefficient Néron--Severi equation and the indefinite generator.\newline \path{returns/S6_02_FAMILY/sections/CD.tex} & New proof\newline §9\\
FAM \nolinkurl{D.1.5}\newline T & Statement~\ref{fam:D.1.5}. Generic rank one and explicit abelian fibres arbitrarily far into the c half-plane.\newline \path{returns/S6_02_FAMILY/sections/CD.tex} & Theorem\newline §9\\
FAM \nolinkurl{D.1.6}\newline T & Statement~\ref{fam:D.1.6}. Exact elliptic-quotient, projective and algebraic-dimension loci.\newline \path{returns/S6_02_FAMILY/sections/CD.tex} & Theorem\newline §9\\
FAM \nolinkurl{D.1.7}\newline T & Statement~\ref{fam:D.1.7}. Actual algebraic-dimension-one fibres, including nontorsion P.\newline \path{returns/S6_02_FAMILY/sections/CD.tex} & Theorem\newline §9\\
FAM \nolinkurl{D.2.1}\newline T & Statement~\ref{fam:D.2.1}. Every compact curve is vertical and every compact surface is a reduced fibre.\newline \path{returns/S6_02_FAMILY/sections/CD.tex} & Theorem\newline §9\\
FAM \nolinkurl{D.2.2}\newline T & Statement~\ref{fam:D.2.2}. Rational and elliptic curves and the singular-fibre descriptions.\newline \path{returns/S6_02_FAMILY/sections/CD.tex} & Theorem\newline §9\\
FAM \nolinkurl{D.2.3}\newline T & Statement~\ref{fam:D.2.3}. All compact curves are classified by the algebraic dimension of their fibre.\newline \path{returns/S6_02_FAMILY/sections/CD.tex} & Theorem\newline §9\\
PUSH D (opening definitions)\newline text & Statement~\ref{push:D.1.setup}. The boundary parameter and its height.\newline \path{returns/S6_02_PUSH/sections/D_push.tex} & Definition\newline §10\\
PUSH \nolinkurl{D.1.1}\newline T & Statement~\ref{push:D.1.1}. The Albanese j-invariant has a natural boundary at |w|=216.\newline \path{returns/S6_02_PUSH/sections/D_push.tex} & Theorem\newline §10\\
PUSH \nolinkurl{D.1.2}\newline T & Statement~\ref{push:D.1.2}. An intrinsic reduced anticanonical surface family exists off a countable exceptional set.\newline \path{returns/S6_02_PUSH/sections/D_push.tex} & Theorem\newline §10\\
PUSH \nolinkurl{D.1.3}\newline T & Statement~\ref{push:D.1.3}. No flat proper analytic extension crosses any point of the outer circle, even fibrewise.\newline \path{returns/S6_02_PUSH/sections/D_push.tex} & Theorem\newline §10\\
PUSH A.1 (opening definitions)\newline text & Statement~\ref{push:A.1.setup}. The ruled normalization, specified pinching Q and linearization line.\newline \path{returns/S6_02_PUSH/sections/AB_push.tex} & Definition\newline §10\\
E26 Conjecture 8.2\newline false as worded & Statement~\ref{e26:Conjecture8.2}. The specified nonnormal pinching is conjectured to fill the untwisted family.\newline \path{papers/arXiv_2609.38442_Engel_src/S6.tex} & Counterexample\newline Refuted by PUSH A.3.2 in text\\
PUSH \nolinkurl{A.3.1}\newline T & Statement~\ref{push:A.3.1}. The intrinsic virtual Euler number equals that of every smooth proper-flat smoothing.\newline \path{returns/S6_02_PUSH/sections/AB_push.tex} & Theorem\newline §10\\
PUSH \nolinkurl{A.3.2}\newline T & Statement~\ref{push:A.3.2}. The specified pinching has virtual Euler number zero and cannot fill Engel’s family.\newline \path{returns/S6_02_PUSH/sections/AB_push.tex} & Counterexample\newline §10\\
PUSH \nolinkurl{A.3.2}\newline A & Statement~\ref{push:A.3.2.replacement}. The nodal and two semistable constructions are precise replacements with different central fibres.\newline \path{returns/S6_02_PUSH/sections/AB_push.tex} & Counterexample\newline Repository:\allowbreak{} PUSH A.1.2 and A.2.2\\
PUSH A.1.1 (proof, cone definition)\newline text & Statement~\ref{push:A.1.fan}. The cones of the two-parameter cusp fan.\newline \path{returns/S6_02_PUSH/sections/AB_push.tex} & Definition\newline §10\\
PUSH \nolinkurl{A.1.1}\newline A & Statement~\ref{push:A.1.1}. The unsubdivided cusp chart is xy=uzv, with cusp coordinate t\_\allowbreak{}c=zv.\newline \path{returns/S6_02_PUSH/sections/AB_push.tex} & Theorem\newline §10; proof in repository:\allowbreak{} PUSH A.1.1\\
PUSH \nolinkurl{A.1.2}\newline A & Statement~\ref{push:A.1.2}. A proper nodal filling contracts the specified (-1,-1) curve on each nonzero fibre.\newline \path{returns/S6_02_PUSH/sections/AB_push.tex} & Theorem\newline §10; proof in repository:\allowbreak{} PUSH A.1.2\\
PUSH A.2.1 (proof, six rays and bottom diagonal)\newline text & Statement~\ref{push:A.2.rays}. The six cusp rays and retained bottom diagonal.\newline \path{returns/S6_02_PUSH/sections/AB_push.tex} & Definition\newline §10\\
PUSH \nolinkurl{A.2.1}\newline A & Statement~\ref{push:A.2.1}. Exactly two subdivisions satisfy the joint-flatness and reducedness restrictions.\newline \path{returns/S6_02_PUSH/sections/AB_push.tex} & Theorem\newline §10; proof in repository:\allowbreak{} PUSH A.2.1\\
PUSH \nolinkurl{A.2.2}\newline A & Statement~\ref{push:A.2.2}. The two natural semistable fillings have the stated normalized central fibres.\newline \path{returns/S6_02_PUSH/sections/AB_push.tex} & Theorem\newline §10; proof in repository:\allowbreak{} PUSH A.2.2\\
PUSH \nolinkurl{A.2.3}\newline A & Statement~\ref{push:A.2.3}. The unrestricted blow-up shape is false; star charts supply its replacement.\newline \path{returns/S6_02_PUSH/sections/AB_push.tex} & Counterexample\newline §10; proof in repository:\allowbreak{} PUSH A.2.3\\
PUSH A.2.4 (proof, base-changed period shears)\newline text & Statement~\ref{push:A.2.periods}. The period shears after the base change u=r\textasciicircum{}n.\newline \path{returns/S6_02_PUSH/sections/AB_push.tex} & Definition\newline §10\\
PUSH \nolinkurl{A.2.4}\newline A & Statement~\ref{push:A.2.4}. Finite fan data enumerate all invariant smooth refinements in the stated scope.\newline \path{returns/S6_02_PUSH/sections/AB_push.tex} & Theorem\newline §10; proof in repository:\allowbreak{} PUSH A.2.4\\
PUSH \nolinkurl{A.2.5}\newline A & Statement~\ref{push:A.2.5}. A terminating blow-up/\allowbreak{}blow-down algorithm describes normalization comparisons.\newline \path{returns/S6_02_PUSH/sections/ab_push_normalization_algorithm.tex} & Counterexample\newline §10; proof in repository:\allowbreak{} PUSH A.2.5\\
PUSH Conventions; B.4.33 (scope)\newline text & Statement~\ref{push:B.4.33.scope}. The fixed sphere family and its distinction from Engel’s untwisted family.\newline \path{returns/S6_02_PUSH/sections/conventions_push.tex} & Definition\newline §11\\
PUSH \nolinkurl{B.4.33}\newline A & Statement~\ref{push:B.4.33}. The retained family has no proper Hausdorff analytic filling after a finite parameter cover.\newline \path{returns/S6_02_PUSH/sections/B_global_noncompletion.tex} & Theorem\newline §11; proof in repository:\allowbreak{} PUSH B.4.33\\
PUSH Conventions (mapping-torus orientation)\newline text & Statement~\ref{push:C.orientation}. Mapping-torus, sphere-gluing and Theta7 orientation conventions.\newline \path{returns/S6_02_PUSH/sections/conventions_push.tex} & Definition\newline §11\\
PUSH C.4.1 (scope)\newline text & Statement~\ref{push:C.4.1.scope}. The exact twist-comparison scope used in PUSH C.3.24.\newline \path{returns/S6_02_PUSH/sections/C_push.tex} & Definition\newline §11\\
PUSH \nolinkurl{C.3.24}\newline A & Statement~\ref{push:C.3.24}. The actual positive-loop mapping class is zero in Theta7, by the orbit-knot method.\newline \path{returns/S6_02_PUSH/sections/C_monodromy_value.tex} & Theorem\newline §11; proof in repository:\allowbreak{} PUSH C.3.24\\
PUSH \nolinkurl{C.3.27}\newline A & Statement~\ref{push:C.3.27}. A spin-c coboundary and direct index additivity give a second zero determination.\newline \path{returns/S6_02_PUSH/sections/C_independent_monodromy_value.tex} & Theorem\newline §11; proof in repository:\allowbreak{} PUSH C.3.27\\
MAPS \nolinkurl{B.6.21}\newline A & Statement~\ref{maps:B.6.21}. The MAPS calculation gives the original (3,4) positive-loop class zero.\newline \path{returns/S6_04_MAPS_MANIFOLDS/sections/maps_B.tex} & Theorem\newline §11; proof in repository:\allowbreak{} MAPS B.6.21\\
PUSH E.1.1 (sphere data); FAM B.3.1 (parameter lattice)\newline text & Statement~\ref{push:E.1.1.scope}. The fixed coprime sphere-pair and twist scope of the higher-family announcements.\newline \path{returns/S6_02_PUSH/sections/E_push.tex}\newline \path{returns/S6_02_FAMILY/sections/AB.tex} & Definition\newline §11\\
MAPS \nolinkurl{B.6.24}\newline A & Statement~\ref{maps:B.6.24}. All stated coprime sphere-pair monodromy classes vanish in Theta7.\newline \path{returns/S6_04_MAPS_MANIFOLDS/sections/maps_B.tex} & Theorem\newline §11; proof in repository:\allowbreak{} MAPS B.6.24\\
PUSH E.2.5 (hypotheses); E.2.7 (parameter)\newline text & Statement~\ref{push:E.2.5.scope}. The precise smooth-specialization hypotheses and fixed-pair parameter.\newline \path{returns/S6_02_PUSH/sections/E_smooth_specialization.tex}\newline \path{returns/S6_02_PUSH/sections/E_zero_period_obstruction.tex} & Definition\newline §11\\
PUSH \nolinkurl{E.2.21}\newline A & Statement~\ref{push:E.2.21}. A smooth completion is an actual member at an interior fixed-pair parameter.\newline \path{returns/S6_02_PUSH/sections/E_global_completion.tex} & Theorem\newline §11; proof in repository:\allowbreak{} PUSH E.2.21\\
PUSH \nolinkurl{E.2.22}\newline A & Statement~\ref{push:E.2.22}. A fixed constructed pair propagates from a nonempty open subset of a disc.\newline \path{returns/S6_02_PUSH/sections/E_disc_propagation.tex} & Theorem\newline §11; proof in repository:\allowbreak{} PUSH E.2.22\\
PUSH \nolinkurl{E.2.23}\newline A & Statement~\ref{push:E.2.23}. Proper smooth deformation chains preserve the ordered coprime multiplicity pair.\newline \path{returns/S6_02_PUSH/sections/E_arbitrary_deformation_bases.tex} & Theorem\newline §11; proof in repository:\allowbreak{} PUSH E.2.23\\
PUSH \nolinkurl{E.2.24}\newline A & Statement~\ref{push:E.2.24}. Each fixed-pair locus is open and closed in the space of integrable structures.\newline \path{returns/S6_02_PUSH/sections/E_continuous_deformations.tex} & Theorem\newline §11; proof in repository:\allowbreak{} PUSH E.2.24\\
\bottomrule\end{longtable}\endgroup
\subsection{Known and new, by section}\label{app:credits}
This table incorporates the source openings of AC26 that the supplied expert table had not performed. \emph{Selected run results} identifies the supplied proof source, not an assertion of absence from all literature.
\begingroup\small\setlength{\tabcolsep}{3pt}\begin{longtable}{@{}L{.05\linewidth}L{.34\linewidth}L{.56\linewidth}@{}}
\toprule § & Known source and locator & Scope and credit\\\midrule\endhead
1 & AC26 Main Theorem; Theorem 3.4; Remark 3.12; DGKXv1 Theorems 7.1--7.2; E26 Section 2. & Construction, canonical bundle, standard intrinsic pencil, factorization and elliptic-logarithm interpretation are known. FAM gives the explicit transported-reference dictionary and corrects E26 Remark 8.1. Its wp-zero geometric reading is credited to AC26.\\
2 & DGKXv1 Theorem 1.4 and bare-period corollary; AC26 Remark 6.4; Viaclovsky Proposition 4.7. & The standard-data 2Z statement is known. The all-discrete-data invariants and two classes at fixed u are the selected FAM results. The 2klZ theorem retains its stated model; the Viaclovsky coordinate dictionary is A.\\
3 & AC26 Theorem 3.4 and Remark 3.12; Eichler--Zagier p.399; Duke--Imamoglu Section 2; Shioda Theorem 8.6. & Finite period values and the elliptic-logarithm interpretation have their source credits. The normalized beta formula, exact nomes, division evaluations and specified integrals are the selected run results. The general division identity is classical local-height quasi-quadraticity, re-derived here; it is not credited as a new general theorem.\\
4 & DGKXv1 Corollary p2:\allowbreak{}cor:\allowbreak{}bare-period; E26 Proposition 2.4. & The abstract punctured disc and sufficiently small geometric parameter are known. The maximum, unique maximizer, exact radius 216 and normalized unit-disc nome are the selected new evaluations. Certificate scope and missing files are explicit.\\
5 & AC26 Proposition 9.13; Theorem 9.1; Remarks 1.3 and 9.21; DGKXv1 Proposition 6.10; Ugarte Theorem 3.2; McHugh Section 7.1; Stelzig Theorem A. & The structure-sheaf direct images, h0q, hp0 and h11=h12+1 are known. The selected cotangent calculation determines h12=1. INTAKE\_\allowbreak{}S6\_\allowbreak{}01 Section 4 also credits the numerical Hodge table to the workbench (GPT-5.6 Sol, Codex reasoning effort ultra), with a degree-count argument; INV supplies the closed scalar and determinant identity. Conditional background formulas for a hypothetical complex S6 are credited where applied; exact Frolicher ranks, evaluated Bott--Chern/\allowbreak{}Aeppli tables and specified cover calculations have their run proofs. Product and metric statements remain A.\\
6 & DGKXv1 Theorem 1.3 and tangent-cohomology theorem; standard proper base change and relative duality. & The tuple (1,1,1,0), Kuranishi line and miniversality are known and re-derived. Integral twists are CDP C5b and INV D4.1. Higher multiplicities carry the two-runs, no-second-reader qualifier. The nearby a=1 corollary is the specified consequence.\\
7 & AC26 Proposition 9.23 and Corollary 9.24; DGKXv1 full-automorphisms theorem. & Aut0 and its action are AC26; the full group is DGKXv1. The conjugate-parameter and real-nome criterion are the selected T results. Fixed loci and classes of real involutions are A.\\
8 & AC26 period monodromies; E26 Remark 8.4. & The construction matrices and Klingler question have source credit. The explicit index-12 subgroup, maximal proper-action open set and period embedding are the selected FAM results.\\
9 & AC26 Lemmas 9.2--9.4, Proposition 9.5, Corollary 9.8, Proposition 4.6; E26 Lemma 2.1; Shioda and Appell--Humbert theory. & The NS equation, independence, generic indefinite class and generic NS rank are already in AC26; the normalization and conductor geometry of W are also there. These results are credited as known content despite the binding T/\allowbreak{}proof placement. The selected torsion-point calculations, all-curve/\allowbreak{}surface statements and special-parameter examples retain their stated scope.\\
10 & E26 Conjecture 8.2; Buchdahl Theorem 11; Grauert base-change and analytic flatness results; analytic specialization sources. & The outer-boundary obstruction and virtual Euler count refute the stated target in their exact source scope. The candidate centre has virtual Euler number zero whereas the smooth sphere fibre has Euler number two. Corrected fillings are A.\\
11 & The supplied PUSH and MAPS returns, with the bounded intake checks. & Zero-end non-filling, trivial Theta7 monodromy and unions of components are A. Two runs use different methods for the monodromy value. The interpretation J=epsilon P.O is an unverified reader remark.\\
12 & INTAKE\_\allowbreak{}FINAL Section 6; RESPONSE\_\allowbreak{}INTAKE; RESPONSE\_\allowbreak{}EXPERT. & The four open question groups are retained as questions, with the MAPS B.6.24 announced scope distinguished from the other families.\\
\bottomrule\end{longtable}\endgroup
\subsection{Dependency concordance}\label{app:dependencies}
The fixed section order retains the source arguments. A pointer to Appendix B is a dependency on its printed re-derivation. Forward references within the selected source order are retained and disclosed in Limits and open points. Every A row points to its repository proof rather than contributing a printed proof.
\begingroup\small\begin{longtable}{@{}L{.21\linewidth}L{.73\linewidth}@{}}
\toprule Result & Uses / dependence\\\midrule\endhead
FAM \nolinkurl{conventions} (\ref{def:conventions}) & AC26 §§2–6; E26 §§2–5\\
FAM \nolinkurl{A.1.1} (\ref{fam:A.1.1}) & Definition~\ref{def:conventions}; \cite[Theorem 7.1]{DGKXv1}; \cite[Theorem 9.1]{AC26}.\\
DGKXv1 Theorem 7.2; thm:\allowbreak{}all-maps-factor (\ref{dgkx:all-maps-factor}) & H²(S⁶;Z)=0; proper image theorem\\
FAM \nolinkurl{B.1.1} (\ref{fam:B.1.1}) & Definition~\ref{def:conventions}; FAM A.3.1 and preceding dictionary identities; \cite[Remark 3.12]{AC26}; \cite[Lemma 2.1; Corollary 2.2; Definition 2.3]{E26}.\\
FAM \nolinkurl{B.1.2} (\ref{fam:B.1.2}) & Definition~\ref{def:conventions}; FAM A.3.1 and preceding dictionary identities; \cite[Remark 3.12]{AC26}; \cite[Lemma 2.1; Corollary 2.2; Definition 2.3]{E26}.\\
FAM \nolinkurl{B.1.3} (\ref{fam:B.1.3}) & Definition~\ref{def:conventions}; FAM A.3.1 and preceding dictionary identities; \cite[Remark 3.12]{AC26}; \cite[Lemma 2.1; Corollary 2.2; Definition 2.3]{E26}.\\
FAM \nolinkurl{B.1.4} (\ref{fam:B.1.4}) & Definition~\ref{def:conventions}; FAM A.3.1 and preceding dictionary identities; \cite[Remark 3.12]{AC26}; \cite[Lemma 2.1; Corollary 2.2; Definition 2.3]{E26}.\\
FAM \nolinkurl{B.2.3} (\ref{fam:B.2.3}) & FAM A.2.1, A.3.1, A.4.1 and B.2.1; \cite[Remark 8.1]{E26}.\\
FAM \nolinkurl{A.5.1} (\ref{fam:A.5.1}) & FAM A.1.1, A.2.1, A.3.1 and A.4.2; \cite[Theorems 1.3, 1.4 and 7.1; punctured-disc corollary]{DGKXv1}.\\
FAM \nolinkurl{B.2.1} (\ref{fam:B.2.1}) & Definition~\ref{def:conventions}; FAM A.1.1--A.3.1 and the preceding discrete-data calculation; \cite[Section 6.4]{AC26}; \cite[Sections 3 and 5]{E26}.\\
FAM \nolinkurl{B.2.2} (\ref{fam:B.2.2}) & Definition~\ref{def:conventions}; FAM A.1.1--A.3.1 and the preceding discrete-data calculation; \cite[Section 6.4]{AC26}; \cite[Sections 3 and 5]{E26}.\\
FAM \nolinkurl{B.3.1} (\ref{fam:B.3.1}) & FAM A.1.2, A.2.1, B.2.1; \cite[Proposition 4.7]{VL26}.\\
PUSH \nolinkurl{E.1.1} (\ref{push:E.1.1}) & FAM B.3.1; PUSH E.1.1--E.1.2; \cite[Section 2]{VL26}.\\
PUSH \nolinkurl{E.1.2} (\ref{push:E.1.2}) & FAM B.3.1; PUSH E.1.1--E.1.2; \cite[Section 2]{VL26}.\\
FAM \nolinkurl{A.1.2} (\ref{fam:A.1.2}) & Definition~\ref{def:conventions}; FAM A.1.1 and preceding appendix calculations.\\
FAM \nolinkurl{A.2.1} (\ref{fam:A.2.1}) & Definition~\ref{def:conventions}; FAM A.1.1 and preceding appendix calculations.\\
FAM \nolinkurl{A.3.1} (\ref{fam:A.3.1}) & Definition~\ref{def:conventions}; FAM A.1.1 and preceding appendix calculations.\\
FAM \nolinkurl{A.4.1} (\ref{fam:A.4.1}) & Definition~\ref{def:conventions}; FAM A.1.1 and preceding appendix calculations.\\
FAM \nolinkurl{A.4.2} (\ref{fam:A.4.2}) & Definition~\ref{def:conventions}; FAM A.1.1 and preceding appendix calculations.\\
WORKBENCH \nolinkurl{pshift:lem:centralizer} (\ref{pshift:lem:centralizer}) & Retained period construction; canonical pencil; finite normalization; explicit centralizer and gluing\\
WORKBENCH \nolinkurl{pshift:thm:classification} (\ref{pshift:thm:classification}) & Retained period construction; canonical pencil; finite normalization; explicit centralizer and gluing\\
FAM \nolinkurl{C.1.1} (\ref{fam:C.1.1}) & FAM B.1.3; the cusp normalization; the Jacobi transformation laws\\
FAM \nolinkurl{C.1.2} (\ref{fam:C.1.2}) & FAM C.1.1; the theta product and its displayed tail bound\\
ARI \nolinkurl{C.1.1} (\ref{ari:C.1.1}) & the theta product and Haar probability measure\\
ARI \nolinkurl{C.1.2} (\ref{ari:C.1.2}) & ARI C.1.1; FAM C.1.1; the fixed cusp branch\\
ARI \nolinkurl{C.2.1} (\ref{ari:C.2.1}) & FAM A.4.1 and A.5.1; ARI C.1.2; the maximum theorem stated in the Result\\
ARI \nolinkurl{C.4.2} (\ref{ari:C.4.2}) & ARI C.1.1; the displayed rational Weierstrass model and standard division-polynomial normalization\\
ARI \nolinkurl{C.4.3} (\ref{ari:C.4.3}) & ARI C.4.2; the nonreal cusp value of the chosen elliptic logarithm\\
ARI \nolinkurl{C.4.4} (\ref{ari:C.4.4}) & ARI C.1.2, C.4.2 and C.4.3; the division-polynomial recurrence\\
ARI \nolinkurl{C.8.6} (\ref{ari:C.8.6}) & ARI C.1.1; the maximum theorem stated in the Result\\
ARI \nolinkurl{C.8.7} (\ref{ari:C.8.7}) & ARI C.8.6; the gamma duplication and triplication formulas\\
ARI \nolinkurl{C.8.8} (\ref{ari:C.8.8}) & the theta series and Poisson summation\\
ARI \nolinkurl{C.8.9} (\ref{ari:C.8.9}) & the normalized theta norm; local root expansions and the cusp branch\\
ARI \nolinkurl{C.8.10} (\ref{ari:C.8.10}) & ARI C.8.9; the height metric and its equidistribution theorem\\
ARI \nolinkurl{C.8.11} (\ref{ari:C.8.11}) & the theta heat equation and the Weierstrass differential equation\\
FAM \nolinkurl{C.2.3} (\ref{fam:C.2.3}) & FAM C.1.1 and C.2.2; the interval cover described below\\
PUSH \nolinkurl{E.4.1} (\ref{push:E.4.1}) & the supplied rational rectangle endpoints; the theta norm; Arb interval arithmetic\\
ARI \nolinkurl{C.3.1} (\ref{ari:C.3.1}) & the theta heat equation; the half-period Weierstrass equation; the cusp-normalized branch\\
ARI \nolinkurl{C.8.1} (\ref{ari:C.8.1}) & the theta product, Poisson summation and Fourier majorants\\
ARI \nolinkurl{C.8.2} (\ref{ari:C.8.2}) & ARI C.8.1 and C.3.1; the theta heat equation and Cauchy estimates\\
ARI \nolinkurl{C.8.3} (\ref{ari:C.8.3}) & ARI C.8.2; the Ramanujan differential identities and Rouche theorem\\
ARI \nolinkurl{C.8.4} (\ref{ari:C.8.4}) & ARI C.4.2, C.4.4 and C.3.1\\
ARI \nolinkurl{C.8.5} (\ref{ari:C.8.5}) & ARI C.8.1–C.8.4; the norm values of ARI C.1.2\\
ARI \nolinkurl{C.8.15} (\ref{ari:C.8.15}) & ARI C.8.11–C.8.14 and the two local root expansions\\
FAM C.1 setup (\nolinkurl{periods:theta-definitions}) & named source setup\\
FAM C.1.2 artifact-availability record (\nolinkurl{periods:special-computation-availability}) & named source setup\\
ARI C.1 setup; C.8 critical setup (\nolinkurl{periods:norm-definition}) & named source setup\\
FAM C.2.1 proof, eq:\allowbreak{}cd-norm (\nolinkurl{fam:eq:cd-norm}) & named source setup\\
ARI C.4 setup and local-height notation in C.4.1 (\nolinkurl{periods:division-model}) & named source setup\\
ARI C.8 measure conventions (\nolinkurl{periods:measure-conventions}) & named source setup\\
FAM;PUSH FAM C.2.3; PUSH E.4.1--3; REVIEW\_\allowbreak{}S6\_\allowbreak{}02\_\allowbreak{}FAMILY §2 (\nolinkurl{periods:certificate-record}) & named source setup\\
ARI C.3 setup (\nolinkurl{periods:local-expansion-constants}) & named source setup\\
ARI C.8 analytic setup (\nolinkurl{periods:ladder-modular-conventions}) & named source setup\\
INV \nolinkurl{notation-cotangent} (\ref{inv:notation-cotangent}) & The family and marking fixed in Section~\ref{sec:family}.\\
INV \nolinkurl{A1.0} (\ref{inv:A1.0}) & The cyclic and cusp charts; analytic base change \cite[III.4.10, III.4.12]{BS76}; analytic duality \cite[pp.~261, 264--266]{RRV71} and \cite[p.~91, Proposition~2]{RR70}.\\
INV \nolinkurl{A1.1} (\ref{inv:A1.1}) & The local cyclic and normal-crossing charts; inv:A1.0.\\
INV \nolinkurl{A1.2} (\ref{inv:A1.2}) & inv:A1.0, inv:A1.1; analytic duality \cite[pp.~261, 264--266]{RRV71}.\\
INV \nolinkurl{notation-period-branch} (\ref{inv:notation-period-branch}) & The period functions and the first-run clockwise convention.\\
INV \nolinkurl{A2.1} (\ref{inv:A2.1}) & inv:A1.0; the period transformations \cite[Definition 3.1 and Theorem 3.4(i)]{AC26}; the modular identities stated in the proof.\\
INV \nolinkurl{A2.2} (\ref{inv:A2.2}) & inv:A2.1; the already identified elliptic section fam:B.1.3; the displayed Ramanujan identities.\\
INV \nolinkurl{A2.4} (\ref{inv:A2.4}) & inv:A1.2, inv:A2.1, inv:A2.2; for the explicitly counterfactual final clause, the residue calculation in inv:A1.3.\\
INV \nolinkurl{A1.3} (\ref{inv:A1.3}) & inv:A1.0--inv:A1.2, inv:A2.1, inv:A2.2, inv:A2.4; the normalized residue sequence; analytic relative duality \cite{RRV71}.\\
INV \nolinkurl{A3.1} (\ref{inv:A3.1}) & inv:A1.0, inv:A1.3; coherent Leray on $\mathbb P^1$ and Serre duality.\\
INV \nolinkurl{notation-hodge-input} (\ref{inv:notation-hodge-input}) & inv:A1.0, inv:A1.3, inv:A3.1; sphere topology.\\
INV \nolinkurl{C1.1} (\ref{inv:C1.1}) & The Hodge table in inv:notation-hodge-input; convergence of the Fr\"olicher spectral sequence to de Rham cohomology.\\
INV \nolinkurl{C2.1} (\ref{inv:C2.1}) & inv:C1.1; the bounded-double-complex decomposition \cite[Theorem 3 and Proposition 6(1)]{Stelzig21}.\\
INV \nolinkurl{C3.3} (\ref{inv:C3.3}) & inv:C1.1, inv:C2.1; repository inv:C1.2 and inv:C3.1; the graded Leibniz rule.\\
INV \nolinkurl{B1.1} (\ref{inv:B1.1}) & The fixed period and cusp charts; inv:A1.0--inv:A1.2; smooth proper base change \cite[III.4.12]{BS76}.\\
INV \nolinkurl{B2.1} (\ref{inv:B2.1-canonical}) & inv:B1.1; \cite[Theorem 1.1 and Proposition 2.3]{VL26}; relative duality \cite{RRV71}.\\
INV \nolinkurl{B1.2} (\ref{inv:B1.2}) & inv:B1.1; the repository residue calculation and relative duality.\\
INV \nolinkurl{B2.1} (\ref{inv:B2.1-hodge}) & inv:B1.2; \cite[Theorem 1.1 and Proposition 2.3]{VL26}.\\
INV \nolinkurl{B3.1} (\ref{inv:B3.1}) & The admissible affine and cusp fillings; inv:B1.1; the elementary covering and normalization calculations in the proof; \cite[III.4.12]{BS76}.\\
INV \nolinkurl{B3.2} (\ref{inv:B3.2}) & inv:B3.1; inv:B1.2 (A, source dependency retained); inv:A1.3; derived relative duality \cite{RRV71}.\\
INV \nolinkurl{B3.3} (\ref{inv:B3.3}) & inv:B3.1, inv:B3.2; the source valuation and monomial argument of inv:B2.2, printed below for $(k,l)=(1,1)$; the marked lattice and affine data.\\
INV \nolinkurl{C3.5} (\ref{inv:C3.5}) & The fibration and sphere topology; Stokes’ theorem.\\
INV \nolinkurl{C3.7} (\ref{inv:C3.7}) & inv:A3.1; the global Dolbeault lift in repository inv:A3.2; inv:C2.1.\\
INV \nolinkurl{A4.1} (\ref{inv:A4.1}) & The presentations of ordinary relative Kähler differentials for $t=xy$, $t=xyz$, and $t=s^m$.\\
INV \nolinkurl{A4.2} (\ref{inv:A4.2}) & inv:A4.1; \cite[Remarks 1.3 and 9.21]{AC26}.\\
FAM \nolinkurl{A.5.2} (\ref{fam:A.5.2}) & The compact family fam:A.5.1 and its nonzero Kodaira--Spencer class fam:A.4.2; \cite[Theorems 1.3 and 6.12]{DGKXv1}.\\
CDP \nolinkurl{notation-integral-twists} (\ref{cdp:notation-integral-twists}) & The standard clockwise affine data and exponential sequence; the root factorization.\\
CDP \nolinkurl{C5b} (\ref{cdp:C5b}) & The root and normalization calculations below; inv:D1.3 for the source-expanded branch-normal and conormal computation; analytic base change \cite[III.4.10, III.4.12]{BS76} and relative duality \cite{RRV71}.\\
PUSH \nolinkurl{E.2.1} (\ref{push:E.2.1}) & The admissible higher-multiplicity data push:E.1; the displayed root-chart and residue computations; \cite[III.4.12]{BS76} and \cite{RRV71}. Two runs, no second reader.\\
FAM F4 consequence of A.5.2; binding row 6d (\ref{fam:A.5.2-corollary}) & fam:A.5.2; \cite[Theorem 9.1(i)]{AC26}.\\
FAM \nolinkurl{B.4.1} (\ref{fam:B.4.1}) & FAM A.1.1, A.1.2, A.2.1, A.3.1 and the cusp valuation comparison in FAM B.2\\
FAM \nolinkurl{B.4.2} (\ref{fam:B.4.2}) & FAM A.3.1, A.5.1, B.2.1 and B.4.1; the cusp-normalized period laws\\
ARI \nolinkurl{C.2.2} (\ref{ari:C.2.2}) & FAM B.4.1, B.4.2 and A.5.1; ARI C.2.1\\
MAPS A.3 (opening conventions) (\ref{maps:A.3.setup}) & FAM B.4.1 and the cusp fan\\
MAPS \nolinkurl{A.3.1} (\ref{maps:A.3.1}) & FAM B.4.1 and B.4.2; the semifree circle action\\
MAPS A.3, equations (A.3.1)--(A.3.2) (\ref{maps:A.3.representatives}) & FAM B.4.2 and the specified root branches\\
MAPS \nolinkurl{A.3.2} (\ref{maps:A.3.2}) & MAPS A.3.1 and its explicit representative data; FAM B.4.2\\
MAPS \nolinkurl{A.3.3} (\ref{maps:A.3.3}) & MAPS A.3.1 and A.3.2\\
FAM F (opening conventions) (\ref{fam:F.setup}) & The retained integral marking and monodromy matrices\\
FAM \nolinkurl{F.1.1} (\ref{fam:F.1.1}) & The matrices $A_1,A_2$ and the integral symplectic form\\
FAM \nolinkurl{F.1.2} (\ref{fam:F.1.2}) & FAM F.1.1 and its exact Heisenberg kernel\\
FAM F.2 (definition preceding F.2.1) (\ref{fam:F.2.setup}) & FAM F.setup\\
FAM \nolinkurl{F.2.1} (\ref{fam:F.2.1}) & FAM F.1.1 and the explicit matrices $J_0,R$\\
FAM \nolinkurl{F.2.2} (\ref{fam:F.2.2}) & FAM F.1.1, F.1.2 and F.2.1; the homogeneous-space stabilizer calculation\\
FAM \nolinkurl{F.2.3} (\ref{fam:F.2.3}) & The global Siegel chart and the signature condition\\
FAM \nolinkurl{F.3.1} (\ref{fam:F.3.1}) & FAM F.2.2 and F.2.3; the finite-root and cusp period expansions\\
FAM \nolinkurl{F.3.2} (\ref{fam:F.3.2}) & The finite special values and the cusp-normalized period matrix\\
FAM \nolinkurl{F.3.3} (\ref{fam:F.3.3}) & The derivative transformation laws and the finite special values\\
FAM \nolinkurl{D.1.1} (\ref{fam:D.1.1}) & The fixed period matrix and the exact division-polynomial computation\\
FAM \nolinkurl{D.1.2} (\ref{fam:D.1.2}) & FAM D.1.1; E26 Lemma 2.1 and its height calculation\\
FAM \nolinkurl{D.1.3} (\ref{fam:D.1.3}) & FAM D.1.1 and the source bound FAM C.2.1\\
FAM D.1.6 (proof, Hermitian-form convention) (\ref{fam:D.1.Hermitian}) & FAM D.1.6, Hermitian-form convention\\
FAM \nolinkurl{D.1.4} (\ref{fam:D.1.4}) & The period matrix; the Hermitian-form convention in FAM D.1.6; AC26 Lemmas 9.2 and 9.4\\
FAM \nolinkurl{D.1.5} (\ref{fam:D.1.5}) & FAM D.1.4; the five-function independence of AC26 Lemma 9.3\\
FAM \nolinkurl{D.1.6} (\ref{fam:D.1.6}) & FAM D.1.4; the Appell--Humbert and Riemann criteria\\
FAM \nolinkurl{D.1.7} (\ref{fam:D.1.7}) & FAM D.1.6 and the cusp expansions\\
FAM \nolinkurl{D.2.1} (\ref{fam:D.2.1}) & $X_c\simeq S^6$; the K\"ahler geometry of the reduced fibres\\
FAM \nolinkurl{D.2.2} (\ref{fam:D.2.2}) & FAM D.1.1 and D.2.1; the finite normalization and boundary identifications of AC26 Proposition 4.6\\
FAM \nolinkurl{D.2.3} (\ref{fam:D.2.3}) & FAM D.1.6 and D.2.1; the Appell--Humbert theorem and its theta equation\\
PUSH D (opening definitions) (\ref{push:D.1.setup}) & The exact radius 216\\
PUSH \nolinkurl{D.1.1} (\ref{push:D.1.1}) & FAM C.1.1 and A.4.1; ARI C.2.1; modular invariance of j\\
PUSH \nolinkurl{D.1.2} (\ref{push:D.1.2}) & Grauert direct-image and proper-mapping theorems with the 1963 correction; Ehresmann; Buchdahl Theorem 11\\
PUSH \nolinkurl{D.1.3} (\ref{push:D.1.3}) & FAM A.1.1; PUSH D.1.1 and D.1.2\\
PUSH A.1 (opening definitions) (\ref{push:A.1.setup}) & E26 Lemma 2.1 and Corollary 2.2\\
E26 Conjecture 8.2 (\ref{e26:Conjecture8.2}) & E26 §8 specified untwisted family\\
PUSH \nolinkurl{A.3.1} (\ref{push:A.3.1}) & The local-complete-intersection deformation calculation; Levy Lemmas 3.1--3.2; \L{}ojasiewicz \S3 Theorems 1--2\\
PUSH \nolinkurl{A.3.2} (\ref{push:A.3.2}) & PUSH A.3.1; the specified pinching; the Euler number two of the chosen untwisted family\\
PUSH \nolinkurl{A.3.2} (\ref{push:A.3.2.replacement}) & PUSH A.1.2 and A.2.2\\
PUSH A.1.1 (proof, cone definition) (\ref{push:A.1.fan}) & The cusp cocharacter lattice\\
PUSH \nolinkurl{A.1.1} (\ref{push:A.1.1}) & PUSH A.1.fan\\
PUSH \nolinkurl{A.1.2} (\ref{push:A.1.2}) & PUSH A.1.1; E26 Lemma 2.1 and Corollary 2.2\\
PUSH A.2.1 (proof, six rays and bottom diagonal) (\ref{push:A.2.rays}) & PUSH A.1.1\\
PUSH \nolinkurl{A.2.1} (\ref{push:A.2.1}) & PUSH A.1.1 and the retained bottom triangulation\\
PUSH \nolinkurl{A.2.2} (\ref{push:A.2.2}) & PUSH A.1.2 and A.2.1\\
PUSH \nolinkurl{A.2.3} (\ref{push:A.2.3}) & PUSH A.2.4 and A.2.5\\
PUSH A.2.4 (proof, base-changed period shears) (\ref{push:A.2.periods}) & PUSH A.2.4, finite fan coordinates\\
PUSH \nolinkurl{A.2.4} (\ref{push:A.2.4}) & PUSH A.1.1 and the explicit base-changed period shears\\
PUSH \nolinkurl{A.2.5} (\ref{push:A.2.5}) & PUSH A.2.4 and its marked star charts\\
PUSH Conventions; B.4.33 (scope) (\ref{push:B.4.33.scope}) & The first-run retained data\\
PUSH \nolinkurl{B.4.33} (\ref{push:B.4.33}) & PUSH B.4.31 and B.4.32; the actual fixed punctured sphere family\\
PUSH Conventions (mapping-torus orientation) (\ref{push:C.orientation}) & The complex orientation and the positive parameter circle\\
PUSH C.4.1 (scope) (\ref{push:C.4.1.scope}) & FAM B.2.1; PUSH C.4.1\\
PUSH \nolinkurl{C.3.24} (\ref{push:C.3.24}) & PUSH C.3.18, C.3.21--C.3.23 and C.4.1\\
PUSH \nolinkurl{C.3.27} (\ref{push:C.3.27}) & PUSH C.3.25 and C.3.26; the shared characteristic calculations C.3.18--C.3.22\\
MAPS \nolinkurl{B.6.21} (\ref{maps:B.6.21}) & MAPS B.6.15--B.6.20 and the specified mapping-torus comparison\\
PUSH E.1.1 (sphere data); FAM B.3.1 (parameter lattice) (\ref{push:E.1.1.scope}) & PUSH E.1.1--E.1.3; FAM B.3.1\\
MAPS \nolinkurl{B.6.24} (\ref{maps:B.6.24}) & The fixed K1 twist datum; MAPS B.6.19 and B.6.22--B.6.23\\
PUSH E.2.5 (hypotheses); E.2.7 (parameter) (\ref{push:E.2.5.scope}) & PUSH E.2.5 and E.2.7\\
PUSH \nolinkurl{E.2.21} (\ref{push:E.2.21}) & PUSH E.2.5, E.2.7 and E.2.10--E.2.20\\
PUSH \nolinkurl{E.2.22} (\ref{push:E.2.22}) & PUSH E.2.1 and E.2.21\\
PUSH \nolinkurl{E.2.23} (\ref{push:E.2.23}) & PUSH E.2.21 and E.2.22; local miniversality\\
PUSH \nolinkurl{E.2.24} (\ref{push:E.2.24}) & PUSH E.2.23 and local completeness of the deformation charts\\
\bottomrule\end{longtable}\endgroup
\subsection{Notation concordance}\label{app:notation}
The standing conventions are those of FAM. The companion notation sheet gives every retained label and the scoped symbol substitutions; the mathematical signs and normalizations are unchanged unless the source explicitly compares two conventions.
\paragraph{root}
FAM conventions are unchanged:\allowbreak{} rho=exp(pi i/\allowbreak{}3); clockwise root deck generators, positive base meridians, Pi\_\allowbreak{}c and w=exp(-pi i c). P\_\allowbreak{}E is renamed P\_\allowbreak{}mark in statements and proofs. Dual covectors u,w appearing in linear combinations and their coefficient comparisons are u\_\allowbreak{}Lambda,w\_\allowbreak{}Lambda; hatted homology vectors retain their names. The shorthand macro for hatted delta is dhat, avoiding the standard T1 command dh. The modular nome is q\_\allowbreak{}tau and the modular discriminant is Delta\_\allowbreak{}mod, agreeing with the periods section; finite discs Delta\_\allowbreak{}j, vector q\_\allowbreak{}n and the auxiliary division-model discriminant keep their distinct meanings.
Scoped renamings avoid collisions with the global parameters u=w\textasciicircum{}2 and w. In FAM A.1.1 the square root of E\_\allowbreak{}6 is h\_\allowbreak{}6, and the two canonical-ring sections u,v are s\_\allowbreak{}can,t\_\allowbreak{}can (the abstract polynomial generators U,V are unchanged). In FAM A.3.1 the preliminary first-coordinate lift is h\_\allowbreak{}loc. In FAM A.4.2 the moving torus coordinate w is zeta\textasciicircum{}new and its projector P is P\_\allowbreak{}KS. In FAM B.1.1, B.1.2 and B.1.4 the line-bundle fibre coordinate W is v, exactly as in the first-run conventions; it is distinct from the cusp fibre W. The reciprocal cusp coordinate v=1/\allowbreak{}t is r\_\allowbreak{}cusp, including the differential in B.1.3. The B.1.3 blow-up parameter u=t-1 is v\_\allowbreak{}t, agreeing with the duplication section. The holomorphic torus-valued unit u(t\_\allowbreak{}c) of B.2.1 is v\_\allowbreak{}unit(t\_\allowbreak{}c). These edits change symbols only, never hypotheses or algebraic identities.
Workbench mathscr V → mathcal V; its homogeneous line mathcal L → mathcal L\_\allowbreak{}per; its local logarithms w\_\allowbreak{}j → ell\_\allowbreak{}j\textasciicircum{}log; its finite base-change variable u → s\_\allowbreak{}new; its local equivariant first-component function u(z) → h\_\allowbreak{}loc(z). The covector coefficient comparison uses u\_\allowbreak{}Lambda,w\_\allowbreak{}Lambda consistently. Both u replacements are confined to the source paragraphs defining those local variables; the global Engel parameter is unchanged. No sign is changed. Its missing labels ffcert:\allowbreak{}thm:\allowbreak{}normalization, prop:\allowbreak{}period-input-audit and thm:\allowbreak{}canonical-ring-recovery are mapped respectively to FAM A.1.2, A.3.1 and A.1.1, whose displayed calculations have been matched.
Every retained FAM st:\allowbreak{}ID label becomes fam:\allowbreak{}ID. The workbench pshift labels are retained. The companion universal map-factorization statement comes from DGKXv1 thm:\allowbreak{}all-maps-factor; it is not part of the supplied FAM A.1.1 statement. DGKX a → c and cA → mathcal P. PUSH E.1.1 and E.1.2 are printed as announcements as mandated, with their statement bodies retained.
\paragraph{periods}
These are literal notation changes in statements and proofs. No change is made to a parameter range, quantifier, equality, inequality or equality locus. Source equation labels receive the namespace fam:\allowbreak{}, ari:\allowbreak{} or push:\allowbreak{}. The retained source statement numbers remain visible.
\begingroup\small\setlength{\tabcolsep}{3pt}\begin{longtable}{@{}L{0.24\linewidth}L{0.25\linewidth}L{0.45\linewidth}@{}}\toprule
Source notation & Paper notation & Meaning and reason\\\midrule\endhead
FAM a=5+2√6 & a\_\allowbreak{}cusp=5+2√6 & Cusp logarithmic constant, distinguished from affine twist variables.\\
FAM C.1 proof u=e\textasciicircum{}(2πiμ₀) & v\_\allowbreak{}cusp & Multiplicative cusp value; avoids Engel's global u=w².\\
FAM C.1 proof scalar f & f\_\allowbreak{}sc & The Weierstrass scaling factor already named f\_\allowbreak{}sc in FAM B.1.1; avoids the threefold projection f.\\
FAM and ARI modular nome q=e\textasciicircum{}(2πiτ) & q\_\allowbreak{}τ & Distinguishes the elliptic modulus from the Albanese nomes.\\
ARI q₁,q₂ & unchanged & q₁=-w²/\allowbreak{}48³, q₂=iw/\allowbreak{}216; q₁=(27/\allowbreak{}64)q₂².\\
ARI η & η\_\allowbreak{}D & Dedekind function; first-run notation, distinguished from the integral cohomology class.\\
ARI N & 𝒩 & The Gaussian-normalized theta norm; FAM C.2 uses 𝒩.\\
ARI L(t)=log N(P\_\allowbreak{}t) & 𝓛(t)=log 𝒩(P\_\allowbreak{}t) & Distinguished from a line bundle. U\_\allowbreak{}Λ retains the full torus Green function.\\
ARI elliptic lattice Λ,Λ′ & Λ\_\allowbreak{}ell,Λ′\_\allowbreak{}ell & Distinguished from FAM's marked rank-four period lattice.\\
ARI Δ₀ & Δ\_\allowbreak{}mod & The modular discriminant q\_\allowbreak{}τ∏(1−q\_\allowbreak{}τⁿ)²⁴.\\
ARI rational model E\_\allowbreak{}t, coordinates x,y, point P & E\_\allowbreak{}t\textasciicircum{}div, x\_\allowbreak{}div,y\_\allowbreak{}div, P\textasciicircum{}div & Auxiliary division-polynomial coordinates, not the FAM cubic coordinates.\\
ARI Δ(t)=27t⁸(t−1)³ & Δ\_\allowbreak{}div(t) & Discriminant of the auxiliary rational model. Generic Δ\_\allowbreak{}E and Δ\_\allowbreak{}Λell mean the discriminant of their named Weierstrass model/\allowbreak{}lattice.\\
ARI C.4.3 radicals W,V & W\_\allowbreak{}rad,V\_\allowbreak{}rad & Products of discriminants and division values, distinguished from the cusp fibre and the integral lattice.\\
ARI C.3.1 G=g₂(i) & G\_\allowbreak{}square & Square-lattice invariant, distinguished from the local function G(h,τ) in the maximum proof.\\
ARI C.3.1 Taylor coefficients A₂,A₄ & α₂,α₄ & Avoids the monodromy matrices A₁,A₂.\\
ARI C.3.1/\allowbreak{}C.8.2 scalar E=π²E₂/\allowbreak{}3 & 𝓔₂ & Scaled quasimodular scalar, distinguished from the rank-one lattice endomorphism E.\\
ARI C.8.2 Fourier sums A\_\allowbreak{}k & 𝒮\_\allowbreak{}k & Scalar theta-series bounds, distinguished from monodromy matrices and duplication polynomials.\\
ARI C.8.4 coefficient bound M & M\_\allowbreak{}coeff & Bound for five polynomial coefficients, distinguished from the global admissibility maximum M.\\
ARI C.8.4 local u=t−1 & v\_\allowbreak{}t & Avoids Engel's multiplicative family parameter.\\
ARI C.8.4 duplication polynomials A\_\allowbreak{}k & 𝒜\_\allowbreak{}k & Avoids the first-run monodromy matrices.\\
ARI C.8.6 integral coordinate u=1/\allowbreak{}℘(z) & v\_\allowbreak{}℘ & Avoids the Engel parameter and the ladder coordinate.\\
ARI C.8.9 local base coordinate w & z\_\allowbreak{}B & Avoids the global moduli parameter w=e\textasciicircum{}(−πic).\\
ARI C.8.11 W & 𝒲 & Logarithmic theta derivative; avoids the cusp fibre W.\\
ARI C.8.11 R=μ′−a & ℛ\_\allowbreak{}μ & Derivative defect; distinguished from radii and the theta quotient.\\
\bottomrule\end{longtable}\endgroup
\paragraph{Cubic normalization and section sign}
FAM B.1.1 uses y\_\allowbreak{}F²=x\_\allowbreak{}F³−3t³(t−1)x\_\allowbreak{}F+2t⁴(t−1)². FAM B.1.4 fixes the cusp branch with distinguished section (0,−√2 t²(t−1)). ARI C.4 uses y\_\allowbreak{}div²=x\_\allowbreak{}div³−(3/\allowbreak{}4)t³(t−1)x\_\allowbreak{}div+(1/\allowbreak{}4)t⁴(t−1)² and section (0,t²(t−1)/\allowbreak{}2). The coordinate substitution x\_\allowbreak{}F=2x\_\allowbreak{}div, y\_\allowbreak{}F=−2√2 y\_\allowbreak{}div identifies both equations and these distinguished sections. This is a coordinate dictionary, not a new geometric theorem. The norm is independent of the Weierstrass homothety and is even in the point, as in the ARI C.4 setup and C.1.1 proof. Generic sigma/\allowbreak{}division-polynomial arguments retain their explicitly stated normalization x=℘, y=℘′/\allowbreak{}2, ψ₁=1, ψ₂=2y.
\paragraph{Special source conventions}
ρ=(1+i√3)/\allowbreak{}2=e\textasciicircum{}(πi/\allowbreak{}3), not the other primitive cube root. The finite values are μ(z₁)=(2−ρ)/\allowbreak{}3, μ(z₂)=(1−i)/\allowbreak{}2; the cusp value is μ₀=1/\allowbreak{}2−i log(5+2√6)/\allowbreak{}(2π). All phases are continued from that branch. The period half-plane is Im c<−M with M=−(3/\allowbreak{}π)log6; the positive-imaginary Albanese representatives are −κ₁ and −κ₂/\allowbreak{}2. The interval cover uses 0≤Re τ≤1/\allowbreak{}2, |τ|≥1 and not a change of root-deck orientation.
The current convention is exactly the corrected ARI C.8.9 convention:\allowbreak{} ddᶜ=(i/\allowbreak{}π)∂\ensuremath{\bar\partial}, positive Betti measure db∧da, total mass 1/\allowbreak{}12, probability measure 12m\_\allowbreak{}P. No sign is changed in the mathematical statements.
\paragraph{cohomology}
The source statements retain the first-run convention:\allowbreak{} the coarse base is B=\textbackslash{}mathbb P\textasciicircum{}1\_\allowbreak{}t, its points are p\_\allowbreak{}1=0, p\_\allowbreak{}2=1, p\_\allowbreak{}0=\textbackslash{}infty, the standard finite multiplicities are 3 and 4, the root generators are clockwise, and \textbackslash{}beta=\textbackslash{}beta\_\allowbreak{}\{\textbackslash{}mathrm\{part\}\}+c. No additional special-point swap is required in these fragments. The alternate Engel marking Q in INV is not used; the paper retains the FAM marking P\_\allowbreak{}\{\textbackslash{}mathrm\{mark\}\}.
\paragraph{Symbol concordance}
\begingroup\small\setlength{\tabcolsep}{3pt}\begin{longtable}{@{}L{0.24\linewidth}L{0.25\linewidth}L{0.45\linewidth}@{}}\toprule
Source symbol & Printed symbol & Reason and scope\\\midrule\endhead
Dual lattice covectors u,w & u\_\allowbreak{}\{\textbackslash{}Lambda\},w\_\allowbreak{}\{\textbackslash{}Lambda\} & Distinguishes the period covectors from family parameters; the hatted lattice vectors stay unchanged. Local module generators named u are untouched.\\
INV logarithmic relative cotangent E or \textbackslash{}mathcal E & \textbackslash{}mathcal E\_\allowbreak{}\{\textbackslash{}rm cot\} & Keeps the relative cotangent bundle separate from the tangent bundle and elliptic curves.\\
FAM/\allowbreak{}PUSH/\allowbreak{}CDP root vertical tangent \textbackslash{}mathcal E & \textbackslash{}mathcal E\_\allowbreak{}\{\textbackslash{}rm tan\} & Distinguishes tangent and cotangent source conventions.\\
INV coarse direct-image bundle G & G\_\allowbreak{}B & Separates the coarse and root-curve extensions.\\
FAM/\allowbreak{}PUSH/\allowbreak{}CDP root direct-image bundle G & G\_\allowbreak{}\{\textbackslash{}mathfrak B\} & Same distinction.\\
INV root curve \textbackslash{}mathcal B & \textbackslash{}mathfrak B & Matches the first-run convention.\\
Root-to-coarse map \textbackslash{}rho & \textbackslash{}varpi & Reserves \textbackslash{}rho for the sixth root e\textasciicircum{}\{\textbackslash{}pi i/\allowbreak{}3\}.\\
INV cotangent determinant line N & N\_\allowbreak{}\{\textbackslash{}rm det\} & Separates this line from the Bott--Chern ideal and from a cusp tube.\\
INV square-zero ideal N & N\_\allowbreak{}\{BC\} & Only the ideal in inv:\allowbreak{}C3.3.\\
FAM/\allowbreak{}PUSH cusp normalization surface E & P\_\allowbreak{}6 & The degree-six del Pezzo surface; not an elliptic curve or a vector bundle.\\
Source B2.2 weighted variables x of weight 4, y of weight 3 & x of weight 3, y of weight 4 & Swapped consistently only in the valuation argument imported into inv:\allowbreak{}B3.3.\\
Source B2.2 auxiliary line \textbackslash{}mathcal L & \textbackslash{}mathcal L\_\allowbreak{}\{\textbackslash{}rm root\} & Separates the valuation argument from the integral Picard generator.\\
CDP residue indices r\_\allowbreak{}1(n),r\_\allowbreak{}2(n) & r\_\allowbreak{}3(n),r\_\allowbreak{}4(n) & The subscripts now record the actual root orders; the residue formulas are unchanged.\\
\bottomrule\end{longtable}\endgroup
Original statement IDs are retained with inv:\allowbreak{}, fam:\allowbreak{}, push:\allowbreak{}, or cdp:\allowbreak{} prefixes. Internal equation labels receive the same run prefix. Label and reference keys are protected before symbol substitutions; no mathematical replacement is allowed inside an identifier. The source's references to Part A now point to sec:\allowbreak{}cotangent-images, except the reference to its Hodge diamond, which points to the more precise sec:\allowbreak{}hodge.
\paragraph{Environments, packages, and macros}
The fragments have no preamble or bibliography. They use theorem, lemma, proposition, corollary, definition, remark, announcement, cited, and proof, with the supplied \textbackslash{}uses\{...\} macro after each formal statement. The central preamble provides these environments and \textbackslash{}uses. Standard requirements are amsmath, amssymb, amsthm, mathtools, mathrsfs, and the central table packages. Existing math operators such as \textbackslash{}Hom, \textbackslash{}Ext, \textbackslash{}rank, \textbackslash{}coker, \textbackslash{}im, \textbackslash{}Spec, \textbackslash{}Proj, \textbackslash{}Pic, \textbackslash{}End, \textbackslash{}Sym, \textbackslash{}Res, \textbackslash{}ord, \textbackslash{}tr, and \textbackslash{}lcm come from that preamble. The proof shorthands use the central \textbackslash{}dbar and \textbackslash{}dhat; the master normalizes any legacy \textbackslash{}dh spelling to \textbackslash{}dhat to avoid redefining the standard T1 text command. No standard LaTeX command is redefined here.
\paragraph{Content-preserving layout edits}
Long proofs are divided into numbered steps by placing step labels at existing source paragraph boundaries. Source formulas are not altered by that operation. References and notation are normalized before recording statement\_\allowbreak{}as\_\allowbreak{}printed; every stored printed statement is matched against the generated fragment. Exact original source statements are retained separately in source\_\allowbreak{}statement.
\paragraph{Final Appendix B cusp-coordinate correction}
The final visual check identified a remaining notation collision in the FAM A.5.2 re-derivation. In that appendix proof, the two occurrences of the weight-twelve modular discriminant \textbackslash{}Delta(\textbackslash{}tau) become \textbackslash{}Delta\_\allowbreak{}\{\textbackslash{}rm mod\}(\textbackslash{}tau). The finite root-disc notation \textbackslash{}Delta\_\allowbreak{}\{s\_\allowbreak{}j\} is unchanged.
The source expressly calls the local equation q=z\_\allowbreak{}0z\_\allowbreak{}1z\_\allowbreak{}2 the original cusp family. Only that coordinate and its subsequent occurrences in the same cusp calculations are renamed to t\_\allowbreak{}c:\allowbreak{} the convergent-series ring, the toric gluing u(q)q\textasciicircum{}\{B\_\allowbreak{}0\textbackslash{}lambda\}, the relation \textbackslash{}Delta(\textbackslash{}tau)=q\textbackslash{},u(q), the field q\textbackslash{}partial\_\allowbreak{}q, and the differentials dq. The last three now read \textbackslash{}Delta\_\allowbreak{}\{\textbackslash{}rm mod\}(\textbackslash{}tau)=t\_\allowbreak{}c\textbackslash{},u(t\_\allowbreak{}c), t\_\allowbreak{}c\textbackslash{}partial\_\allowbreak{}\{t\_\allowbreak{}c\}, and dt\_\allowbreak{}c. Cohomological indices q in higher direct images, degree conditions and cases are unchanged. No equality between t\_\allowbreak{}c and the modular nome q\_\allowbreak{}\textbackslash{}tau is introduced.
The exact scope and replacement counts are recorded in \path{appendix_notation_changes.json}; appendix\_\allowbreak{}notation\_\allowbreak{}fix.diff records the before/\allowbreak{}after appendix fragment. The generating script applies the same finite list of replacements with count assertions, preserving all statements, hypotheses and proof steps.
The same final scan found exactly one additional modular-discriminant token in the main INV A1.0 proof. With the assembling editor's explicit authorization, \textbackslash{}Delta(\textbackslash{}tau)\textasciicircum{}\{-1\} there is likewise renamed to \textbackslash{}Delta\_\allowbreak{}\{\textbackslash{}rm mod\}(\textbackslash{}tau)\textasciicircum{}\{-1\}. Its surrounding cusp coordinate was already t\_\allowbreak{}c; no other main-fragment notation is changed. This additional one-token correction is included in \path{appendix_notation_changes.json}.
\paragraph{geometry ends}
All mathematical conventions are those of FAM, \path{returns/S6_02_FAMILY/sections/conventions.tex}. The changes below are notation, reference or layout changes; the source statements, hypotheses and scopes remain separately recorded in \path{records.json}. Exact statement and proof diffs are in \path{statement_comparisons.json} and \path{proof_comparisons.json}.
\paragraph{Common conventions retained}
The base is B=P\textasciicircum{}1\_\allowbreak{}t, with p\_\allowbreak{}1=0, p\_\allowbreak{}2=1, p\_\allowbreak{}0=infinity; the retained multiplicities are (3,4). Homology columns have order (hat-gamma,hat-u,hat-w,hat-delta). The deck generators are clockwise root transformations, g\_\allowbreak{}1g\_\allowbreak{}2g\_\allowbreak{}0=1; positive geometric parameter loops are counterclockwise. The period matrix is Pi\_\allowbreak{}c=((6mu,tau,1,0),(beta\_\allowbreak{}part+c,mu,0,1)), with (beta\_\allowbreak{}part+tau)(p\_\allowbreak{}0)=0, and w=exp(-pi*i*c). The Albanese nome is q\_\allowbreak{}2=iw/\allowbreak{}216. The fibre vector is e=(0,1)\textasciicircum{}t=Pi\_\allowbreak{}c hat-delta; the lattice vector is d=hat-delta.
\paragraph{Symbol mappings}
\begingroup\small\setlength{\tabcolsep}{3pt}\begin{longtable}{@{}L{0.17\linewidth}L{0.21\linewidth}L{0.22\linewidth}L{0.34\linewidth}@{}}\toprule
Source and scope & Source notation & Printed notation & Reason /\allowbreak{} sign convention\\\midrule\endhead
FAM F, lattice & e=hat-delta & d=hat-delta & Distinguishes the primitive lattice vector from the fibre vector e already used in FAM AB. Every concatenated occurrence Je, C e, R e, Z e, stabilizer subscript and q\_\allowbreak{}e is changed consistently.\\
FAM F & E=e tensor gamma, q\_\allowbreak{}e, P\_\allowbreak{}\{e,Z/\allowbreak{}R/\allowbreak{}Q\} & E=d tensor gamma, q\_\allowbreak{}d, P\_\allowbreak{}\{d,Z/\allowbreak{}R/\allowbreak{}Q\} & Same operator, invariant function and stabilizer.\\
FAM F, auxiliary scalars & matrix entry/\allowbreak{}denominator d & d\_\allowbreak{}aux & Prevents the local scalar from colliding with the renamed lattice vector. Matrix formulas and denominator-clearing integers are otherwise unchanged.\\
FAM F.3.3 & cusp q=exp(2pi*i*s) & t\_\allowbreak{}c=exp(2pi*i*s) & Keeps the cusp coordinate distinct from the Albanese nome. Exponentials e\textasciicircum{}(pi*i/\allowbreak{}3) and e\textasciicircum{}(2pi*i*s) remain exponentials.\\
FAM F and D, dual coframe & u,w in wedge forms & u\_\allowbreak{}Lambda,w\_\allowbreak{}Lambda & Matches the root paper's dual-coframe notation; hatted homology vectors and parameter w are unchanged.\\
FAM D & beta & beta=beta\_\allowbreak{}part+c & Retains the actual period, not the particular solution alone.\\
FAM D.1.4 & T,U,D in G & T=Im tau, U=Im mu, D=Im beta-6U\textasciicircum{}2/\allowbreak{}T stated explicitly & T is defined in FAM C; U is not locally defined in \path{CD.tex}. This notational ambiguity is recorded in the discrepancy report. No entry of G, determinant or intersection sign is changed.\\
FAM D.1.4/\allowbreak{}D.1.6 & Hermitian form & H\_\allowbreak{}E(v,w)=Etilde(v,Jw)-i Etilde(v,w) & The source convention, linear in its first entry, is retained. Its Chern form is -Im H\_\allowbreak{}E; no sign is silently imported from AC26.\\
MAPS A.3 & homology g & hat-gamma in the translation formula & The short alias is removed.\\
MAPS A.3 representative formula & S(t)d in a fibre-coordinate translation, although d=hat-delta & S(t)e, with e=Pi\_\allowbreak{}c d=(0,1)\textasciicircum{}t explicitly stated & Exposes the source's lattice/\allowbreak{}fibre-vector shorthand. All coefficients, branches and real representatives are unchanged.\\
MAPS A.3 & circle action rho & rho\_\allowbreak{}circ(b) is translation by b e, b in R/\allowbreak{}Z & Explicitly identifies the source circle action with FAM B.4.1 and reserves the unadorned rho=exp(pi*i/\allowbreak{}3) for the period constant. Every rho(1/\allowbreak{}2) is renamed consistently.\\
MAPS A.3 & D\_\allowbreak{}1,D\_\allowbreak{}2,D\_\allowbreak{}12 & unchanged & Directions e\_\allowbreak{}1,e\_\allowbreak{}2,e\_\allowbreak{}1-e\_\allowbreak{}2; D\_\allowbreak{}2 is the circle-fixed curve. These are not interchanged.\\
PUSH A & Engel line bundle L & M\_\allowbreak{}S=O\_\allowbreak{}S(P-O) & Matches FAM conventions. The cotangent complex mathbb L\_\allowbreak{}D in A.2.5 remains mathbb L\_\allowbreak{}D; it is not the line bundle.\\
PUSH A & linearization line J & J\_\allowbreak{}lin & Avoids confusing the line bundle with a real complex structure or the monodromy integer. Its Chern class is still the lower-case j.\\
PUSH A fan & cusp q, height h\_\allowbreak{}q & t\_\allowbreak{}c, h\_\allowbreak{}\{t\_\allowbreak{}c\} & Same cusp base coordinate and same primitive ray heights. The degeneration parameter u is unchanged.\\
E26 Conjecture 8.2 quotation & M & M inside the verbatim quotation, identified with M\_\allowbreak{}S immediately afterward & Preserves the exact conjecture wording; outside that quotation M remains the period maximum.\\
PUSH E.2.7/\allowbreak{}E.2.21 & fixed-pair parameter z=exp(-pi*i*c/\allowbreak{}(kl)) & w\_\allowbreak{}\{k,l\} & Separates the parameter from the first-run uniformizing coordinate z.\\
MAPS B.6.24 & positive generator 2kl of the parameter-period lattice & unchanged & The source positive closing-translation convention is retained. The positive geometric w loop in the retained family lifts to c -> c-2; no nonzero sign-sensitive mapping class is inferred.\\
FAM/\allowbreak{}PUSH bibliography & DGKX26v1, DuGuoKusak\allowbreak{}abeXie2026\allowbreak{}v1 & DGKXv1 & Common bibliography key. The full-automorphism locator is the theorem labelled p2:\allowbreak{}thm:\allowbreak{}full-automorphisms, not an unqualified period-classification citation.\\
FAM period cross-reference & eq:\allowbreak{}cd-beta & fam:\allowbreak{}eq:\allowbreak{}cd-beta & Matches the period agent's sole equation label; no second label is inserted into its equation.\\
\bottomrule\end{longtable}\endgroup
\paragraph{Statement and equation labels}
Every retained numbered source statement has the label run:\allowbreak{}original-ID:\allowbreak{} fam:\allowbreak{}B.4.1, fam:\allowbreak{}B.4.2, ari:\allowbreak{}C.2.2, maps:\allowbreak{}A.3.1--maps:\allowbreak{}A.3.3, fam:\allowbreak{}F.1.1--fam:\allowbreak{}F.3.3, fam:\allowbreak{}D.1.1--fam:\allowbreak{}D.2.3, the retained push:\allowbreak{}A.*, push:\allowbreak{}D.1.*, and all specified push:\allowbreak{}B.4.33, push:\allowbreak{}C.3.24, push:\allowbreak{}C.3.27, maps:\allowbreak{}B.6.21, maps:\allowbreak{}B.6.24, push:\allowbreak{}E.2.21--push:\allowbreak{}E.2.24.
The exact E26 target is labelled e26:\allowbreak{}Conjecture8.2. The final replacement sentence of PUSH A.3.2 is labelled push:\allowbreak{}A.3.2.replacement and separately typed A; its source ID remains A.3.2. The rest of that source statement is the T-item push:\allowbreak{}A.3.2.
FAM displayed labels acquire the fam:\allowbreak{} prefix, including fam:\allowbreak{}eq:\allowbreak{}F-g, fam:\allowbreak{}eq:\allowbreak{}F-Gamma, fam:\allowbreak{}eq:\allowbreak{}F-law, fam:\allowbreak{}eq:\allowbreak{}F-SLpresentation, fam:\allowbreak{}eq:\allowbreak{}F-H, fam:\allowbreak{}eq:\allowbreak{}F-qD, fam:\allowbreak{}eq:\allowbreak{}F-tauaction, fam:\allowbreak{}eq:\allowbreak{}F-j, fam:\allowbreak{}eq:\allowbreak{}cd-NS and fam:\allowbreak{}eq:\allowbreak{}cd-theta-sections. Contextual definitions use labels ending .setup, .scope, .representatives, .fan, .rays, .periods, .Hermitian or .orientation; their metadata identifies the original unnumbered context or proof extract explicitly, rather than inventing an original theorem number.
\paragraph{Scoped reuse retained from the sources}
The symbols B (base /\allowbreak{} local 2-by-2 block), T (modular generator /\allowbreak{} imaginary period /\allowbreak{} translation t\_\allowbreak{}\{6P\}), and E (rank-one marking operator /\allowbreak{} alternating class) have explicitly separated local scopes, as in the first run. The local S(t) in the real representative table is a scalar translation function, whereas S in the filling section is Engel's rational elliptic surface. The record does not identify these objects.
\subsection{Verification record, one row per check}\label{app:verification}
R: statement read without following its proof. F: proof or explicitly delimited proof portion followed. D: steps re-derived. C: numerical or exact recomputation. S: source opened and matched. No L check is recorded. A producing-run computation is identified as such. The historical records do not constitute a fresh check of this assembly.
\begingroup\footnotesize\setlength{\tabcolsep}{3pt}\begin{longtable}{@{}L{.08\linewidth}L{.31\linewidth}L{.56\linewidth}@{}}
\toprule Check & By / date / type / scope & Verdict, findings, independence and report\\\midrule\endhead
V001 & Claude Fable 5.2, effort high\newline 2026-11-12; F\newline FAM A.1.1 & correct. Characters chi\_\allowbreak{}j(0) = zeta\_\allowbreak{}j\textasciicircum{}(1+a\_\allowbreak{}j), a = (2,1), from det R\_\allowbreak{}j at tau = rho, i; orders (2,1,-1) of Theta = E4\textasciicircum{}2 sqrt(E6)/\allowbreak{}Delta; divisor of Theta\textasciicircum{}-1 e dt = 2S2 - W, hence K = O(-2S2), equal to [AC26] Theorem (4) (T l. 156); dimension count floor(m/\allowbreak{}2)+1 for C[U,V]. K\_\allowbreak{}X formula is in [AC26]; ring and direct image are the workbench's (F17).\newline Given the return, bibliography and programs; one AI reader.\newline \path{status/REVIEW_S6_02_FAMILY.md}, Section 3, A.1.1\\
V002 & Claude Fable 5.2, effort high\newline 2026-11-12; F,C\newline FAM A.1.2 & correct. Fixed and moving lattices and their saturation by program; degrees 9, 8 from the isogeny kernels of orders 3, 2. Standard (workbench FF1-FF46).\newline Given the return, bibliography and programs; one AI reader.\newline \path{status/REVIEW_S6_02_FAMILY.md}, Section 3, A.1.2\\
V003 & Claude Fable 5.2, effort high\newline 2026-11-12; F,D,C\newline FAM A.2.1 & correct. Whole argument by hand; centralizer QI + QE, N1, N2, quotients Z/\allowbreak{}3 and Z/\allowbreak{}4 + Z/\allowbreak{}2 by my program from T1, T2 of [AC26]; the deck equation and its m\_\allowbreak{}j-fold iterate.\newline Given the return, bibliography and programs; one AI reader.\newline \path{status/REVIEW_S6_02_FAMILY.md}, Section 3, A.2.1\\
V004 & Claude Fable 5.2, effort high\newline 2026-11-12; F,C\newline FAM A.3.1 & correct. (I-A\_\allowbreak{}j)r\_\allowbreak{}j identities, l\_\allowbreak{}0 = (3n/\allowbreak{}2)w - (n/\allowbreak{}2)d, N\_\allowbreak{}j l\_\allowbreak{}j by program; first-coordinate sheaf = O(-p0) from the orders of Theta; H\textasciicircum{}1(V) = 0, H\textasciicircum{}0(V) = C e; invariance of the A2 fan under integer shears; equality of cusp multipliers at c and c+n. Properness at the cusp:\allowbreak{} read.\newline Given the return, bibliography and programs; one AI reader.\newline \path{status/REVIEW_S6_02_FAMILY.md}, Section 3, A.3.1\\
V005 & Claude Fable 5.2, effort high\newline 2026-11-12; F\newline FAM A.4.1 & correct. Albanese lattices; Im(kappa1\textasciicircum{}part - kappa2\textasciicircum{}part) = -0.13735807... by my own evaluation (inside the certified (-0.2,-0.1)).\newline Given the return, bibliography and programs; one AI reader.\newline \path{status/REVIEW_S6_02_FAMILY.md}, Section 3, A.4.1\\
V006 & Claude Fable 5.2, effort high\newline 2026-11-12; F\newline FAM A.4.2 & correct. P\textasciicircum{}2 = P, derivative P/\allowbreak{}(2iD); transfer to the threefold through the bi-anticanonical pencil:\allowbreak{} read. Consistent with A.5.2, where the generator of H\textasciicircum{}1(T) is e tensor [gamma].\newline Given the return, bibliography and programs; one AI reader.\newline \path{status/REVIEW_S6_02_FAMILY.md}, Section 3, A.4.2\\
V007 & Claude Fable 5.2, effort high\newline 2026-11-12; F\newline FAM A.5.1 & known. Assembly checked; the review identifies Theorem 1.4 of DGKXv1 and notes its absence in v2. No date comparison is carried into the paper.\newline Given the return, bibliography and programs; one AI reader.\newline \path{status/REVIEW_S6_02_FAMILY.md}, Section 3, A.5.1\\
V008 & Claude Fable 5.2, effort high\newline 2026-11-12; F,D\newline FAM A.5.2 & known. Every term recomputed:\allowbreak{} Pic of the root curve P(3,4), K = O(-7), L = O(-1), E = O + O(-1), R\textasciicircum{}q f\_\allowbreak{}* O = O, O + O(-1), O(-1); acyclicity of O(k), k = 1,-1,-2,-5,-6; h(T(-log W)) = (1,1,0,0); quotient nu\_\allowbreak{}* K\_\allowbreak{}dP6; result (1,1,1,0). Independent check (mine):\allowbreak{} Hirzebruch-Riemann-Roch gives chi(T\_\allowbreak{}X) = c3/\allowbreak{}2 = 1. Stated in arXiv:\allowbreak{}2610.00250v1 (Theorem 1.3 and introduction); absent from v2.\newline Given the return, bibliography and programs; one AI reader.\newline \path{status/REVIEW_S6_02_FAMILY.md}, Section 3, A.5.2\\
V009 & Claude Fable 5.2, effort high\newline 2026-11-12; F,C\newline FAM B.1.1 & correct. Line bundle O(P-O), factors (1, e\textasciicircum{}(2 pi i mu)); discriminant -108 t\textasciicircum{}8 (t-1)\textasciicircum{}3 and j = 1728t by program. MW = Z P with height 1/\allowbreak{}6 is the Oguiso-Shioda entry for IV*, III, I1 (used by Engel, Lemma 2.1).\newline Given the return, bibliography and programs; one AI reader.\newline \path{status/REVIEW_S6_02_FAMILY.md}, Section 3, B.1.1\\
V010 & Claude Fable 5.2, effort high\newline 2026-11-12; F,C\newline FAM B.1.2 & correct. det of the real period matrix = T D (numerical check).\newline Given the return, bibliography and programs; one AI reader.\newline \path{status/REVIEW_S6_02_FAMILY.md}, Section 3, B.1.2\\
V011 & Claude Fable 5.2, effort high\newline 2026-11-12; F,C\newline FAM B.1.3 & correct. Height 2 - 4/\allowbreak{}3 - 1/\allowbreak{}2 = 1/\allowbreak{}6; wp((1-i)/\allowbreak{}2, i) = wp((2-rho)/\allowbreak{}3, rho) = 0 numerically. The return's correction of the course label (P\_\allowbreak{}* meets E\_\allowbreak{}(2,-), not E\_\allowbreak{}(2,+)) is right:\allowbreak{} Y2 = y/\allowbreak{}t\textasciicircum{}2 = sqrt2 (t-1) -> -sqrt2; course l. 680 says E\_\allowbreak{}(2,+).\newline Given the return, bibliography and programs; one AI reader.\newline \path{status/REVIEW_S6_02_FAMILY.md}, Section 3, B.1.3\\
V012 & Claude Fable 5.2, effort high\newline 2026-11-12; F\newline FAM B.1.4 & correct. div of the transported section = F\_\allowbreak{}infinity:\allowbreak{} read. lambda0 = 1 holds because the reference linearization is defined to be the transported one; see 5.4.\newline Given the return, bibliography and programs; one AI reader.\newline \path{status/REVIEW_S6_02_FAMILY.md}, Section 3, B.1.4\\
V013 & Claude Fable 5.2, effort high\newline 2026-11-12; F,C\newline FAM B.2.1 & correct. Norm criterion against (B2.2) by brute force, 604,800 data, no mismatch; p = -r1 mod 3, p = r2 mod 4; cusp necessity (B2.3) and sufficiency:\allowbreak{} read.\newline Given the return, bibliography and programs; one AI reader.\newline \path{status/REVIEW_S6_02_FAMILY.md}, Section 3, B.2.1\\
V014 & Claude Fable 5.2, effort high\newline 2026-11-12; F\newline FAM B.2.2 & correct. Two square roots of u; signs from a\_\allowbreak{}(j-1) = -P\_\allowbreak{}E v\_\allowbreak{}j/\allowbreak{}m\_\allowbreak{}j and gamma(v1) = 1, gamma(v2) = -1. P\_\allowbreak{}E conjugates A\_\allowbreak{}j to Engel's T\_\allowbreak{}0 (Thm 5.1) and T\_\allowbreak{}1 ((6.1), display garbled in my text copy).\newline Given the return, bibliography and programs; one AI reader.\newline \path{status/REVIEW_S6_02_FAMILY.md}, Section 3, B.2.2\\
V015 & Claude Fable 5.2, effort high\newline 2026-11-12; F\newline FAM B.2.3 & correct. Engel's Remark 8.1 read:\allowbreak{} "a proper, holomorphic submersion Y -> Delta\textasciicircum{}* whose fiber over u is X(u, a0, a1)". By Construction 3.5 and Cor. 3.9, psi(a0) is in (1/\allowbreak{}3)Z minus Z, so the datum is not invariant under the u-monodromy. Second witness:\allowbreak{} j(Alb S2) has a w\textasciicircum{}-1 term.\newline Given the return, bibliography and programs; one AI reader.\newline \path{status/REVIEW_S6_02_FAMILY.md}, Section 3, B.2.3\\
V016 & Claude Fable 5.2, effort high\newline 2026-11-12; F,D\newline FAM B.3.1 & correct with a gap. Arithmetic (N\_\allowbreak{}3k = k N1, N\_\allowbreak{}4l = l N2, lcm(k,2l) = 2kl, CRT) checked by hand. Gap:\allowbreak{} the map from Viaclovsky's data (m\_\allowbreak{}i, a\_\allowbreak{}i), v\_\allowbreak{}i = r\_\allowbreak{}i delta + n\_\allowbreak{}i u\_\allowbreak{}i, to (r\_\allowbreak{}j, s\_\allowbreak{}j, l\_\allowbreak{}0) is not written; sufficiency is "as in B.2.1".\newline Given the return, bibliography and programs; one AI reader.\newline \path{status/REVIEW_S6_02_FAMILY.md}, Section 3, B.3.1\\
V017 & Claude Fable 5.2, effort high\newline 2026-11-12; F,C\newline FAM B.4.1 & known. Coboundary lattice has determinant 1 in the allowed pairs (program); H\textasciicircum{}0(V) = C e. "Aut(X\_\allowbreak{}a) = C\textasciicircum{}*" is in Theorem 1.4 of arXiv:\allowbreak{}2610.00250v1.\newline Given the return, bibliography and programs; one AI reader.\newline \path{status/REVIEW_S6_02_FAMILY.md}, Section 3, B.4.1\\
V018 & Claude Fable 5.2, effort high\newline 2026-11-12; F,C\newline FAM B.4.2 & correct. H\textasciicircum{}2 = I, H A2 H = A2\textasciicircum{}-1, H A1 H = A2\textasciicircum{}-1 A1\textasciicircum{}-1 A2, Pi\_\allowbreak{}(-cbar) H = -conj(Pi), twisted targets (-eps + d, eps' - 3d) by program. Independent check (mine):\allowbreak{} with Re kappa2\textasciicircum{}part = 3/\allowbreak{}2 and Re kappa1\textasciicircum{}part = 1/\allowbreak{}2, the Albanese moduli of the conjugate manifold match those of X\_\allowbreak{}(c') exactly for c' = 1 - cbar mod 2 and not for c' = -cbar.\newline Given the return, bibliography and programs; one AI reader.\newline \path{status/REVIEW_S6_02_FAMILY.md}, Section 3, B.4.2\\
V019 & Claude Fable 5.2, effort high\newline 2026-11-12; F,C\newline FAM C.1.1 & correct. The two Jacobi laws give (beta1) of [AC26] exactly (hand, and numerically with the explicit branch); cusp limit from sin\textasciicircum{}2(pi mu0) = 3; R(z2) = e\textasciicircum{}(pi/\allowbreak{}4) sqrt2 to 40 digits.\newline Given the return, bibliography and programs; one AI reader.\newline \path{status/REVIEW_S6_02_FAMILY.md}, Section 3, C.1.1\\
V020 & Claude Fable 5.2, effort high\newline 2026-11-12; F,C\newline FAM C.1.2 & correct. All five values reproduced to 40 digits by different code.\newline Given the return, bibliography and programs; one AI reader.\newline \path{status/REVIEW_S6_02_FAMILY.md}, Section 3, C.1.2\\
V021 & Claude Fable 5.2, effort high\newline 2026-11-12; F,D,C\newline FAM C.2.3 & correct. Section 2:\allowbreak{} full re-run of the interval proof. Formulas and tails of `wp\_\allowbreak{}iv`, `N2\_\allowbreak{}iv` checked against the text; half-period identity log N = Re H - pi (Im h)\textasciicircum{}2/\allowbreak{}T checked; cusp bound 1.0814 < sqrt2 by hand. The rigorous part rests on mpmath's interval arithmetic; no second interval library was used.\newline Given the return, bibliography and programs; one AI reader.\newline \path{status/REVIEW_S6_02_FAMILY.md}, Section 3, C.2.3\\
V022 & Claude Fable 5.2, effort high\newline 2026-11-12; F\newline FAM D.1.1 & correct. x(6P) has denominator 648 (11t\textasciicircum{}2 - 6016t + 8192)\textasciicircum{}2 (my own chord-tangent computation); roots 64(47 +- 27 sqrt3)/\allowbreak{}11.\newline Given the return, bibliography and programs; one AI reader.\newline \path{status/REVIEW_S6_02_FAMILY.md}, Section 3, D.1.1\\
V023 & Claude Fable 5.2, effort high\newline 2026-11-12; F\newline FAM D.1.2 & known. Abelian at torsion points:\allowbreak{} index-N\textasciicircum{}3 product sublattice checked; 3k\textasciicircum{}2 - 1 from Engel's Lemma 2.1 (read in the text).\newline Given the return, bibliography and programs; one AI reader.\newline \path{status/REVIEW_S6_02_FAMILY.md}, Section 3, D.1.2\\
V024 & Claude Fable 5.2, effort high\newline 2026-11-12; F,D,C\newline FAM D.1.3 & correct. At t = 32/\allowbreak{}5, y(2P) = 0 (exact order 4); L-images (-101 tau, 2 tau + 2, 3, 4); c = -2.25 - 34.0210972421\allowbreak{}8209472308 i reproduced to 20 digits; Im c < -22 by hand. Scope in 5.3.\newline Given the return, bibliography and programs; one AI reader.\newline \path{status/REVIEW_S6_02_FAMILY.md}, Section 3, D.1.3\\
V025 & Claude Fable 5.2, effort high\newline 2026-11-12; F\newline FAM D.1.4 & correct. NS equation from the wedge of the rows of Pi; det = T D < 0; eta\textasciicircum{}2 = -12 in the complex orientation.\newline Given the return, bibliography and programs; one AI reader.\newline \path{status/REVIEW_S6_02_FAMILY.md}, Section 3, D.1.4\\
V026 & Claude Fable 5.2, effort high\newline 2026-11-12; F\newline FAM D.1.5 & correct. Pf = -3n\textasciicircum{}2 for the displayed class; independence argument read.\newline Given the return, bibliography and programs; one AI reader.\newline \path{status/REVIEW_S6_02_FAMILY.md}, Section 3, D.1.5\\
V027 & Claude Fable 5.2, effort high\newline 2026-11-12; F\newline FAM D.1.6 & known. Standard theory of 2-tori (Appell-Humbert, Riemann conditions); the cited lecture notes not opened by me.\newline Given the return, bibliography and programs; one AI reader.\newline \path{status/REVIEW_S6_02_FAMILY.md}, Section 3, D.1.6\\
V028 & Claude Fable 5.2, effort high\newline 2026-11-12; F,D\newline FAM D.1.7 & correct. Pf(a eta + b E) = 6a\textasciicircum{}2 by hand.\newline Given the return, bibliography and programs; one AI reader.\newline \path{status/REVIEW_S6_02_FAMILY.md}, Section 3, D.1.7\\
V029 & Claude Fable 5.2, effort high\newline 2026-11-12; F\newline FAM D.2.1 & correct. Uses H\textasciicircum{}2(X;Z) = 0; argument complete.\newline Given the return, bibliography and programs; one AI reader.\newline \path{status/REVIEW_S6_02_FAMILY.md}, Section 3, D.2.1\\
V030 & Claude Fable 5.2, effort high\newline 2026-11-12; F\newline FAM D.2.2 & correct. Read.\newline Given the return, bibliography and programs; one AI reader.\newline \path{status/REVIEW_S6_02_FAMILY.md}, Section 3, D.2.2\\
V031 & Claude Fable 5.2, effort high\newline 2026-11-12; F\newline FAM D.2.3 & known. A parametrization (theta divisors) more than a list; standard; the cited notes not opened by me.\newline Given the return, bibliography and programs; one AI reader.\newline \path{status/REVIEW_S6_02_FAMILY.md}, Section 3, D.2.3\\
V032 & Claude Fable 5.2, effort high\newline 2026-11-12; F,D,C\newline FAM F.1.1 & correct. By program:\allowbreak{} Omega-invariance, A1 = g(a,(1,-1),-2), A2 = g(S,(0,-1),3), alpha(S,T,a) = (9,1,4), chi a homomorphism (300 random products), chi = 0 on 300 random words, K, U, V, [U,V] = Z\textasciicircum{}-1, chi(I+E) = 1. The index argument by hand.\newline Given the return, bibliography and programs; one AI reader.\newline \path{status/REVIEW_S6_02_FAMILY.md}, Section 3, F.1.1\\
V033 & Claude Fable 5.2, effort high\newline 2026-11-12; F\newline FAM F.1.2 & correct. Read; follows from the group law.\newline Given the return, bibliography and programs; one AI reader.\newline \path{status/REVIEW_S6_02_FAMILY.md}, Section 3, F.1.2\\
V034 & Claude Fable 5.2, effort high\newline 2026-11-12; F,C\newline FAM F.2.1 & correct. J0\textasciicircum{}2 = -I, J0\textasciicircum{}t Omega J0 = Omega, eigenvalues (1,-1,6,-6), R = g(T\textasciicircum{}72, 0, -12) in Gamma, R J0 = J0 R, (R-I)\textasciicircum{}2 = 0 by program.\newline Given the return, bibliography and programs; one AI reader.\newline \path{status/REVIEW_S6_02_FAMILY.md}, Section 3, F.2.1\\
V035 & Claude Fable 5.2, effort high\newline 2026-11-12; F,D\newline FAM F.2.2 & correct. Transitivity with compact stabilizer on each sign region; density of rational conjugates of J0 in the null locus; the 12th-power step. Redone by hand.\newline Given the return, bibliography and programs; one AI reader.\newline \path{status/REVIEW_S6_02_FAMILY.md}, Section 3, F.2.2\\
V036 & Claude Fable 5.2, effort high\newline 2026-11-12; F,C\newline FAM F.2.3 & correct. q\_\allowbreak{}e = -6/\allowbreak{}D checked numerically at a random period point.\newline Given the return, bibliography and programs; one AI reader.\newline \path{status/REVIEW_S6_02_FAMILY.md}, Section 3, F.2.3\\
V037 & Claude Fable 5.2, effort high\newline 2026-11-12; F\newline FAM F.3.1 & correct. j(tau) = 1728 t is a holomorphic left inverse.\newline Given the return, bibliography and programs; one AI reader.\newline \path{status/REVIEW_S6_02_FAMILY.md}, Section 3, F.3.1\\
V038 & Claude Fable 5.2, effort high\newline 2026-11-12; F\newline FAM F.3.2 & correct. Read.\newline Given the return, bibliography and programs; one AI reader.\newline \path{status/REVIEW_S6_02_FAMILY.md}, Section 3, F.3.2\\
V039 & Claude Fable 5.2, effort high\newline 2026-11-12; F,D\newline FAM F.3.3 & correct. mu' = -tau'/\allowbreak{}3, beta' = 2 tau'/\allowbreak{}3 at z1 recomputed by hand.\newline Given the return, bibliography and programs; one AI reader.\newline \path{status/REVIEW_S6_02_FAMILY.md}, Section 3, F.3.3\\
V040 & Claude Fable 5.2, effort high\newline 2026-11-12; C\newline ab\_\allowbreak{}exact & identical. Python 3.13.9; sympy 1.13.1; mpmath 1.3.0. The review identifies the complete outputs as identical.\newline Rerun of the shipped program in a scratch copy.\newline \path{status/REVIEW_S6_02_FAMILY.md}, Section 2\\
V041 & Claude Fable 5.2, effort high\newline 2026-11-12; C\newline ab\_\allowbreak{}geometry & identical. Python 3.13.9; sympy 1.13.1; mpmath 1.3.0. The review identifies the complete outputs as identical.\newline Rerun of the shipped program in a scratch copy.\newline \path{status/REVIEW_S6_02_FAMILY.md}, Section 2\\
V042 & Claude Fable 5.2, effort high\newline 2026-11-12; C\newline ab\_\allowbreak{}deformation & identical. Python 3.13.9; sympy 1.13.1; mpmath 1.3.0. The review identifies the complete outputs as identical.\newline Rerun of the shipped program in a scratch copy.\newline \path{status/REVIEW_S6_02_FAMILY.md}, Section 2\\
V043 & Claude Fable 5.2, effort high\newline 2026-11-12; C\newline cd\_\allowbreak{}bounds & identical. Python 3.13.9; sympy 1.13.1; mpmath 1.3.0. The review identifies the complete outputs as identical.\newline Rerun of the shipped program in a scratch copy.\newline \path{status/REVIEW_S6_02_FAMILY.md}, Section 2\\
V044 & Claude Fable 5.2, effort high\newline 2026-11-12; C\newline cd\_\allowbreak{}exact & identical. Python 3.13.9; sympy 1.13.1; mpmath 1.3.0. The review identifies the complete outputs as identical.\newline Rerun of the shipped program in a scratch copy.\newline \path{status/REVIEW_S6_02_FAMILY.md}, Section 2\\
V045 & Claude Fable 5.2, effort high\newline 2026-11-12; C\newline cd\_\allowbreak{}local\_\allowbreak{}bound & identical. Python 3.13.9; sympy 1.13.1; mpmath 1.3.0. The review identifies the complete outputs as identical.\newline Rerun of the shipped program in a scratch copy.\newline \path{status/REVIEW_S6_02_FAMILY.md}, Section 2\\
V046 & Claude Fable 5.2, effort high\newline 2026-11-12; C\newline cd\_\allowbreak{}multiples\_\allowbreak{}initial & identical. Python 3.13.9; sympy 1.13.1; mpmath 1.3.0. The review identifies the complete outputs as identical.\newline Rerun of the shipped program in a scratch copy.\newline \path{status/REVIEW_S6_02_FAMILY.md}, Section 2\\
V047 & Claude Fable 5.2, effort high\newline 2026-11-12; C\newline cd\_\allowbreak{}periods & identical. Python 3.13.9; sympy 1.13.1; mpmath 1.3.0. The review identifies the complete outputs as identical.\newline Rerun of the shipped program in a scratch copy.\newline \path{status/REVIEW_S6_02_FAMILY.md}, Section 2\\
V048 & Claude Fable 5.2, effort high\newline 2026-11-12; C\newline cd\_\allowbreak{}special\_\allowbreak{}intervals & identical. Python 3.13.9; sympy 1.13.1; mpmath 1.3.0. The review identifies the complete outputs as identical.\newline Rerun of the shipped program in a scratch copy.\newline \path{status/REVIEW_S6_02_FAMILY.md}, Section 2\\
V049 & Claude Fable 5.2, effort high\newline 2026-11-12; C\newline cd\_\allowbreak{}torsion\_\allowbreak{}independent & identical. Python 3.13.9; sympy 1.13.1; mpmath 1.3.0. The review identifies the complete outputs as identical.\newline Rerun of the shipped program in a scratch copy.\newline \path{status/REVIEW_S6_02_FAMILY.md}, Section 2\\
V050 & Claude Fable 5.2, effort high\newline 2026-11-12; C\newline f\_\allowbreak{}exact & identical. Python 3.13.9; sympy 1.13.1; mpmath 1.3.0. The review identifies the complete outputs as identical.\newline Rerun of the shipped program in a scratch copy.\newline \path{status/REVIEW_S6_02_FAMILY.md}, Section 2\\
V051 & Claude Fable 5.2, effort high\newline 2026-11-12; C\newline cd\_\allowbreak{}numerics & identical apart from last platform digits. Two double-precision grid cross-checks differed only in their last digits.\newline Methods exactly as recorded by the review.\newline \path{status/REVIEW_S6_02_FAMILY.md}, Section 2\\
V052 & Claude Fable 5.2, effort high\newline 2026-11-12; C\newline f\_\allowbreak{}interval\_\allowbreak{}verify & identical. 19,911 leaves; exact area 23/\allowbreak{}40; 19,701 norm boxes; 36 local boxes; 174 outside the unit circle.\newline Methods exactly as recorded by the review.\newline \path{status/REVIEW_S6_02_FAMILY.md}, Section 2\\
V053 & Claude Fable 5.2, effort high\newline 2026-11-12; C\newline all 19,701 norm boxes & all recertified. Own \path{recertify_all.py} runner called the supplied certify routine; identical upper-bound strings; 73 seconds, 12 processes. Same mpmath interval library, not a new implementation.\newline Methods exactly as recorded by the review.\newline \path{status/REVIEW_S6_02_FAMILY.md}, Section 2\\
V054 & Claude Fable 5.2, effort high\newline 2026-11-12; C\newline independent exact lattice calculations & checked. \path{v1_exact.py} used the AC26 matrices for FAM A.2, A.3, B.2, B.4, F.1 and F.2.\newline Methods exactly as recorded by the review.\newline \path{status/REVIEW_S6_02_FAMILY.md}, Section 2\\
V055 & Claude Fable 5.2, effort high\newline 2026-11-12; C\newline independent period numerics & checked as numerical evidence. \path{v2_numeric.py}:\allowbreak{} 6,651 grid points, the modular arc and circles around i; independent mpmath theta evaluation.\newline Methods exactly as recorded by the review.\newline \path{status/REVIEW_S6_02_FAMILY.md}, Section 2\\
V056 & Claude Fable 5.2, effort high\newline 2026-11-12; C\newline independent elliptic arithmetic & checked. \path{v3_ec.py} recomputed the D.1 arithmetic and C.2.2 constants.\newline Methods exactly as recorded by the review.\newline \path{status/REVIEW_S6_02_FAMILY.md}, Section 2\\
V057 & Claude Fable 5.2, effort high\newline 2026-11-12; C\newline independent special values & checked as numerical evidence. \path{v4_closed_forms.py} evaluated the closed forms listed in Section 4.9; proofs were not supplied by that numerical check.\newline Methods exactly as recorded by the review.\newline \path{status/REVIEW_S6_02_FAMILY.md}, Section 2\\
V058 & Claude Opus 5.5, effort max\newline 2026-11-12; S\newline DGKXv1 three theorem statements and v2 contents & known labels confirmed. The Kuranishi, exact-period and tangent-cohomology statements were matched in v1; their absence in v2 was checked.\newline Both supplied TeX sources opened by the main session.\newline \path{status/REVIEW_S6_02_FAMILY.md}, added main-session note\\
V059 & Claude Fable 5.2, effort max\newline 2026-11-16; R,S\newline DGKXv1 and v2 scope & matched statements; proofs not read. Abstract and introduction; 17 statements in Sections 6--7; three Appendix C propositions; Appendix D opening; full term counts. For v2:\allowbreak{} abstract, introduction statements and section list.\newline Expert had the TeX and the status discussion.\newline \path{status/RESPONSE_EXPERT.md}, Section 1\\
V060 & Claude Fable 5.2, effort max\newline 2026-11-13, recorded 2026-11-16; D,C\newline ARI C.8.6 and C.8.7 & derived and numerically tested. Mahler-measure identity and two beta integrals derived by hand and tested to seven digits. This is not a claim of a fresh line-by-line read of the source proof.\newline Expert derivations reported in the supplied response.\newline \path{status/RESPONSE_EXPERT.md}, Section 4\\
V061 & Claude Fable 5.2, effort high\newline 2026-11-15; D\newline ARI C.4.2 & checks agree. Hand checks at t=1, 32/\allowbreak{}5 and 0; no code rerun.\newline Bounded arithmetic intake and its recorded hand checks.\newline \path{status/INTAKE_S6_05_FINAL.md}, Section 1; \path{INTAKE_FINAL.md}, Section 4\\
V062 & Claude Fable 5.2, effort high\newline 2026-11-15; F,D\newline ARI C.8.5 displayed inequalities & displayed inequalities agree. Every displayed inequality and the size of the three uniform bounds were checked. C.8.1--C.8.3 were not read line by line. Code and PDF were not read.\newline Given the source proof and intake notes.\newline \path{status/INTAKE_S6_05_FINAL.md}, Section 1\\
V063 & GPT-6 Astra, ChatGPT Pro mode\newline PUSH producing run; C\newline PUSH E.4 Arb certification & all boxes certified as reported. 19,701 boxes; largest upper enclosure below 1.999877571767410525. The intake reports this result and did not rerun the code. Detailed precision is in the source certificate block.\newline Separate interval library; same rational box endpoints.\newline \path{status/INTAKE_FINAL.md}, Section 1, item 4; PUSH E.4\\
V064 & Claude Fable 5.2, effort high\newline 2026-11-16; F\newline INV A1.3; direct images & no error found in the portion read. Proof read; A1.1--A1.2, which supply local freeness and direct images, were not read.\newline Bounded intake; precise read limits retained.\newline \path{status/INTAKE_S6_01.md}, opening and Section 2\\
V065 & Claude Fable 5.2, effort high\newline 2026-11-16; F\newline INV A2.1; period determinant identity & no error found in the portion read. Statement and full proof read.\newline Bounded intake; precise read limits retained.\newline \path{status/INTAKE_S6_01.md}, opening and Section 2\\
V066 & Claude Fable 5.2, effort high\newline 2026-11-16; F\newline INV A2.2; the fourth-root scalar & no error found in the portion read. Statement and full proof read; separate scalar computation is listed below.\newline Bounded intake; precise read limits retained.\newline \path{status/INTAKE_S6_01.md}, opening and Section 2\\
V067 & Claude Fable 5.2, effort high\newline 2026-11-16; F\newline INV A2.4; nonzero Kodaira--Spencer scalar & no error found in the portion read. Statement and full proof read.\newline Bounded intake; precise read limits retained.\newline \path{status/INTAKE_S6_01.md}, opening and Section 2\\
V068 & Claude Fable 5.2, effort high\newline 2026-11-16; F,D\newline INV A3.1; Hodge table & no error found in the portion read. Table recomputed from the six direct-image sheaves; Serre duality, Euler characteristic and the AC26 relation checked.\newline Bounded intake; precise read limits retained.\newline \path{status/INTAKE_S6_01.md}, opening and Section 2\\
V069 & Claude Fable 5.2, effort high\newline 2026-11-16; F\newline INV C1.1; Frolicher sequence & no error found in the portion read. Every arrow followed; no higher differential has both nonzero source and target; middle rows exact.\newline Bounded intake; precise read limits retained.\newline \path{status/INTAKE_S6_01.md}, opening and Section 2\\
V070 & Claude Fable 5.2, effort high\newline 2026-11-16; F,D\newline INV C2.1; Bott--Chern and Aeppli numbers & no error found in the portion read. First half of proof read; numbers recounted from the double-complex arrows. Remaining proof was not fully read.\newline Bounded intake; precise read limits retained.\newline \path{status/INTAKE_S6_01.md}, opening and Section 2\\
V071 & Claude Fable 5.2, effort high\newline 2026-11-16; F\newline INV C3.3; square-zero product & no error found in the portion read. Proof read and Leibniz argument followed; it uses C3.1, which was not read. Binding status remains A.\newline Bounded intake; precise read limits retained.\newline \path{status/INTAKE_S6_01.md}, opening and Section 2\\
V072 & Claude Fable 5.2, effort high\newline 2026-11-16; F,D\newline INV B2.1; higher-multiplicity canonical divisor & no error found in the portion read. Proof read; the two orders and the fibre degree checked. Higher-multiplicity Hodge inputs were not read.\newline Bounded intake; precise read limits retained.\newline \path{status/INTAKE_S6_01.md}, opening and Section 2\\
V073 & Claude Fable 5.2, effort high\newline 2026-11-16; F\newline INV B3.2; degree-n covers & no error found in the portion read. Proof read; the cusp residue mechanism and Betti agreement identified. B3.1 was not read.\newline Bounded intake; precise read limits retained.\newline \path{status/INTAKE_S6_01.md}, opening and Section 2\\
V074 & Claude Fable 5.2, effort high\newline 2026-11-16; F\newline INV B3.3; cover cohomology & no error found in the portion read. Statement and first half of proof read; not a complete proof check.\newline Bounded intake; precise read limits retained.\newline \path{status/INTAKE_S6_01.md}, opening and Section 2\\
V075 & Claude Fable 5.2, effort high\newline 2026-11-16; C\newline INV A2.2 scalar & 40 digits agree. \path{s601_scalar_check.py} reproduced 0.2366093140\allowbreak{}8078522468\allowbreak{}7705113450\allowbreak{}4496062943 and j second derivative -162 Gamma(1/\allowbreak{}4)\textasciicircum{}8/\allowbreak{}pi\textasciicircum{}4.\newline Independent numerical route using the local period expansion and j second derivative.\newline \path{status/INTAKE_S6_01.md}, Section 2\\
V076 & Claude Fable 5.2, effort high\newline 2026-11-16; F,D\newline PUSH A.3.1 & no gap found in the portion read. Full proof read; virtual Euler count redone as 24-(12+12)+0=0.\newline Bounded intake with earlier notes and source text.\newline \path{status/INTAKE_FINAL.md}, opening and Section 1\\
V077 & Claude Fable 5.2, effort high\newline 2026-11-16; F\newline PUSH A.3.2 & no gap found in the portion read. Full proof read; the specified candidate cannot smooth to the sphere fibre.\newline Bounded intake with earlier notes and source text.\newline \path{status/INTAKE_FINAL.md}, opening and Section 1\\
V078 & Claude Fable 5.2, effort high\newline 2026-11-16; F\newline PUSH D.1.1 & no gap found in the portion read. Full proof read; Albanese modular invariant natural boundary argument.\newline Bounded intake with earlier notes and source text.\newline \path{status/INTAKE_FINAL.md}, opening and Section 1\\
V079 & Claude Fable 5.2, effort high\newline 2026-11-16; F\newline PUSH D.1.2 & no gap found in the portion read. Full proof read; base change on the curve and Buchdahl input checked in the argument.\newline Bounded intake with earlier notes and source text.\newline \path{status/INTAKE_FINAL.md}, opening and Section 1\\
V080 & Claude Fable 5.2, effort high\newline 2026-11-16; R\newline PUSH D.1.3 & no gap found in the portion read. Opening read, not its whole proof.\newline Bounded intake with earlier notes and source text.\newline \path{status/INTAKE_FINAL.md}, opening and Section 1\\
V081 & Claude Fable 5.2, effort high\newline 2026-11-16; F\newline PUSH B.4.4 & no gap found in the portion read. Marked-map extension proof read.\newline Bounded intake with earlier notes and source text.\newline \path{status/INTAKE_FINAL.md}, opening and Section 1\\
V082 & Claude Fable 5.2, effort high\newline 2026-11-16; F\newline PUSH B.4.27 & no gap found in the portion read. Log-torus sheaf cohomology proof read.\newline Bounded intake with earlier notes and source text.\newline \path{status/INTAKE_FINAL.md}, opening and Section 1\\
V083 & Claude Fable 5.2, effort high\newline 2026-11-16; F\newline PUSH B.4.30 & no gap found in the portion read. Base modification and extension-of-differences proof read.\newline Bounded intake with earlier notes and source text.\newline \path{status/INTAKE_FINAL.md}, opening and Section 1\\
V084 & Claude Fable 5.2, effort high\newline 2026-11-16; F\newline PUSH B.4.31 & no gap found in the portion read. Global-meridian proof read.\newline Bounded intake with earlier notes and source text.\newline \path{status/INTAKE_FINAL.md}, opening and Section 1\\
V085 & Claude Fable 5.2, effort high\newline 2026-11-16; F\newline PUSH B.4.32 & no gap found in the portion read. Global-loop defect proof read.\newline Bounded intake with earlier notes and source text.\newline \path{status/INTAKE_FINAL.md}, opening and Section 1\\
V086 & Claude Fable 5.2, effort high\newline 2026-11-16; F\newline PUSH B.4.33 & no gap found in the portion read. Proof read; explicit objects B.4.9, B.4.18 and B.4.19 were not independently re-derived.\newline Bounded intake with earlier notes and source text.\newline \path{status/INTAKE_FINAL.md}, opening and Section 1\\
V087 & Claude Fable 5.2, effort high\newline 2026-11-16; F\newline PUSH C.3.21 & no gap found in the portion read. Actual delta-square proof read.\newline Bounded intake with earlier notes and source text.\newline \path{status/INTAKE_FINAL.md}, opening and Section 1\\
V088 & Claude Fable 5.2, effort high\newline 2026-11-16; F\newline PUSH C.3.22 & no gap found in the portion read. Signed-cap comparison proof read.\newline Bounded intake with earlier notes and source text.\newline \path{status/INTAKE_FINAL.md}, opening and Section 1\\
V089 & Claude Fable 5.2, effort high\newline 2026-11-16; F\newline PUSH C.3.23 & no gap found in the portion read. Equivariant integer proof read.\newline Bounded intake with earlier notes and source text.\newline \path{status/INTAKE_FINAL.md}, opening and Section 1\\
V090 & Claude Fable 5.2, effort high\newline 2026-11-16; F\newline PUSH C.3.24 & no gap found in the portion read. Monodromy-value proof read; C.3.10, C.3.14 and C.3.18 were not read in this intake.\newline Bounded intake with earlier notes and source text.\newline \path{status/INTAKE_FINAL.md}, opening and Section 1\\
V091 & Claude Fable 5.2, effort high\newline 2026-11-16; F\newline PUSH C.3.27 & no gap found in the portion read. Second determination proof read.\newline Bounded intake with earlier notes and source text.\newline \path{status/INTAKE_FINAL.md}, opening and Section 1\\
V092 & Claude Fable 5.2, effort high\newline 2026-11-16; F\newline PUSH E.2.21 & no gap found in the portion read. Disc-completion proof read.\newline Bounded intake with earlier notes and source text.\newline \path{status/INTAKE_FINAL.md}, opening and Section 1\\
V093 & Claude Fable 5.2, effort high\newline 2026-11-16; F\newline PUSH E.2.22 & no gap found in the portion read. Connected-base propagation proof read; E.2.5--E.2.20 remain unread inputs.\newline Bounded intake with earlier notes and source text.\newline \path{status/INTAKE_FINAL.md}, opening and Section 1\\
V094 & Claude Fable 5.2, effort high\newline 2026-11-16; F\newline PUSH E.2.24 & no gap found in the portion read. Continuous-deformation proof read; the Kuranishi theorem is an external input.\newline Bounded intake with earlier notes and source text.\newline \path{status/INTAKE_FINAL.md}, opening and Section 1\\
V095 & Claude Fable 5.2, effort high\newline 2026-11-16; F\newline MAPS A.3.1 & no conflict found. Full proof read. A.3.2 was read only as statement and setup; fixed-locus conclusions retain A status.\newline Bounded intake.\newline \path{status/INTAKE_FINAL.md}, opening and Sections 3,5\\
V096 & Claude Fable 5.2, effort high\newline 2026-11-16; F\newline ARI C.2.2 & agrees with the real-parameter criterion. Full proof read; compared with FAM B.4.2 and MAPS A.3.\newline Bounded intake.\newline \path{status/INTAKE_FINAL.md}, opening and Section 5\\
V097 & Claude Fable 5.2, effort high\newline 2026-11-16; D\newline CDP C5b at n=-1,0,12 & Euler characteristic agrees. The supplied response records the three values; INTAKE\_\allowbreak{}FINAL gives the n=0 tuple (1,1,1,0).\newline Hand substitution, not a full independent proof.\newline \path{status/RESPONSE_INTAKE.md}, Section 2; \path{INTAKE_FINAL.md}, Section 2\\
V098 & Claude Fable 5.2, effort high\newline 2026-11-16; R\newline PUSH E.2.1 and CDP E11 & same higher-multiplicity tangent tuple. Two runs, no second reader of these proofs; this qualifier remains attached to the T statement.\newline Cross-run comparison of separate producing sessions.\newline \path{status/INTAKE_FINAL.md}, Section 5 X3; \path{RESPONSE_INTAKE.md}, Section 2\\
V099 & Claude Opus 5.5, effort max\newline Recorded 2026-11-16; D\newline special norm and phase calculation & n=3 division case checked by intake. The intake records NOTES Section 33(a) as a valid derivation of the special norm and phases, and Section 34 as the exact nome-ratio calculation. No unseen notes are asserted to have been opened here.\newline Main-session calculation as evaluated by the intake.\newline \path{status/INTAKE_FINAL.md}, Section 5, NOTES Sections 33--34 discussion\\
V100 & GPT-6 Astra, ChatGPT Pro mode\newline 2026-11-16; S\newline AC26; Main Theorem items (1)--(5); Remark 3.12; Remark 6.4; Theorem 9.1; Section 10 & opened and matched. Construction on S6, algebraic dimension one, canonical bundle; elliptic-logarithm generator; classification left open in Remark 6.4; the source of the counterexample in the prescribed introduction paragraph. Checked in the supplied AI transcription, not in an original PDF.\newline Assembly with supplied source and review context; not an independent proof review.\newline \path{records/root/citation_checks.json}; \path{original/AC26_transcription.tex}\\
V101 & GPT-6 Astra, ChatGPT Pro mode\newline 2026-11-16; S\newline E26; Section 2; Lemma 2.1; Corollary 2.2; Definition 2.3; Proposition 2.4; Sections 3,5,6; Remark 8.1 and Conjecture 8.2 & opened and matched. Pointed cubic, Mordell-Weil generator and reference lift up to scale; inverse-translation convention; small-parameter compactness; lattice basis and affine data; exact correction and counterexample targets.\newline Assembly with supplied source and review context; not an independent proof review.\newline \path{records/root/citation_checks.json}; \path{papers/arXiv_2609.38442_Engel_src/S6.tex}\\
V102 & GPT-6 Astra, ChatGPT Pro mode\newline 2026-11-16; S\newline DGKXv1; Theorems 1.3,1.4; Theorem 7.1 and thm:\allowbreak{}all-maps-factor; p2:\allowbreak{}thm:\allowbreak{}full-automorphisms; p2:\allowbreak{}cor:\allowbreak{}bare-period; Appendix C & opened and matched. Standard-data exact 2Z identification, tangent tuple and Kuranishi germ, intrinsic bi-anticanonical pencil, unique factorization of every map to P1, full automorphism group and unevaluated punctured-disc radius. All-affine-data scope of FAM A.1.1 is not identical to the cited standard-data theorem.\newline Assembly with supplied source and review context; not an independent proof review.\newline \path{records/root/citation_checks.json}; \path{papers/arXiv_2610.00250v1_Du_Guo_Kusakabe_Xie.tex}\\
V103 & GPT-6 Astra, ChatGPT Pro mode\newline 2026-11-16; S\newline DGKXv2; Abstract, introduction, section input list; full-text term searches & opened and matched. Oka content retained; the v1 exact-period, tangent-cohomology, intrinsic-pencil and Kuranishi sections and statements are absent. No reason for this omission is inferred.\newline Assembly with supplied source and review context; not an independent proof review.\newline \path{records/root/citation_checks.json}; \path{papers/arXiv_2610.00250v2_Du_Guo_Kusakabe_Xie.tex}\\
V104 & GPT-6 Astra, ChatGPT Pro mode\newline 2026-11-16; S\newline VL26; Theorem 1.1; Proposition 4.7, label nonbih, and proof & opened and matched. The source's higher multiplicities and non-isotriviality argument; the factorization of maps to the base used for the stated-model proof. The full coordinate dictionary remains an announcement.\newline Assembly with supplied source and review context; not an independent proof review.\newline \path{records/root/citation_checks.json}; \path{papers/arXiv_2610.39102_Viaclovsky_log_transformations_src/Logarithmic_S6_October_28_2026_Final.tex}\\
V105 & GPT-6 Astra, ChatGPT Pro mode\newline 2026-11-16; S\newline OS91; Main Theorem table entry 49; Section 3 equation (3.2) & opened and matched. Root lattice E6 plus A1, determinant six, no torsion and Mordell-Weil lattice with generator height 1/\allowbreak{}6; height input for the IV*,III,I1 surface.\newline Assembly with supplied source and review context; not an independent proof review.\newline \path{records/root/citation_checks.json}; \url{https://rikkyo.repo.nii.ac.jp/record/10006/files/AA00610867_40-01_08.pdf}\\
V106 & GPT-6 Astra, ChatGPT Pro mode\newline 2026-11-16; S\newline Shi90; Theorem 8.6, equations (8.11)--(8.12), page 228; local contribution table (8.16), page 229 & opened and matched. Height pairing formula and local contributions 1/\allowbreak{}2 at III and 4/\allowbreak{}3 at IV*. The source-run citation to Table 8.9 is an incorrect locator; the matched table is (8.16).\newline Assembly with supplied source and review context; not an independent proof review.\newline \path{records/root/citation_checks.json}; \url{https://rikkyo.repo.nii.ac.jp/record/10027/files/AA00610867_39-02_05.pdf}\\
V107 & GPT-6 Astra, ChatGPT Pro mode\newline 2026-11-16; S\newline CDP20; Corollary 2.3 & opened and matched. The asserted algebraic dimension zero for a compact complex threefold homeomorphic to S6; used solely in the prescribed introduction paragraph.\newline Assembly with supplied source and review context; not an independent proof review.\newline \path{records/root/citation_checks.json}; \url{https://www-fourier.univ-grenoble-alpes.fr/~demailly/source_files/cdp/cdp-corrigendum-v5.pdf}\\
V108 & GPT-6 Astra, ChatGPT Pro mode\newline 2026-11-16; S\newline AC26; Theorem 3.4(ii)--(iv); Remark 3.12 & opened and matched. Finite mu values, period laws, invariant D and elliptic-logarithm interpretation; none of the theta norm, exact radius, integral, division or measure claims found by the recorded searches.\newline Assembly with supplied source and review context; not an independent proof review.\newline \path{records/periods/citation_checks.json}; \path{original/AC26_transcription.tex}\\
V109 & GPT-6 Astra, ChatGPT Pro mode\newline 2026-11-16; S\newline E26; Relation to the original manuscript; §2; Remark 8.1 & opened and matched. Geometric interpretation of the four periods and variation of the elliptic curves with linearization parameter; exact nomes and constants not present.\newline Assembly with supplied source and review context; not an independent proof review.\newline \path{records/periods/citation_checks.json}; \path{papers/arXiv_2609.38442_Engel_src/S6.tex}\\
V110 & GPT-6 Astra, ChatGPT Pro mode\newline 2026-11-16; S\newline DGKXv1; Theorem 1.4 (p2:\allowbreak{}thm:\allowbreak{}exact-period); Corollary p2:\allowbreak{}cor:\allowbreak{}bare-period; definition of M\_\allowbreak{}0 & opened and matched. Theorem 1.4 states bare and pair isomorphism iff parameter difference lies in 2Z; Corollary p2:\allowbreak{}cor:\allowbreak{}bare-period gives 0<|w|<exp(-pi M\_\allowbreak{}0); M\_\allowbreak{}0 is not evaluated.\newline Assembly with supplied source and review context; not an independent proof review.\newline \path{records/periods/citation_checks.json}; \path{papers/arXiv_2610.00250v1_Du_Guo_Kusakabe_Xie.tex}\\
V111 & GPT-6 Astra, ChatGPT Pro mode\newline 2026-11-16; S\newline EZ82; Unnumbered theorem p.399 and Fourier corollary p.400 & opened and matched. The zeros of Weierstrass wp are the elliptic integral branches with cusp constant log(5+2sqrt6)/\allowbreak{}(2pi i).\newline Assembly with supplied source and review context; not an independent proof review.\newline \path{records/periods/citation_checks.json}; \url{https://people.mpim-bonn.mpg.de/zagier/files/doi/10.1007/BF01453974/fulltext.pdf}\\
V112 & GPT-6 Astra, ChatGPT Pro mode\newline 2026-11-16; S\newline DI10; §2, equation (6), unnumbered Eichler--Zagier theorem; Theorem 1 & opened and matched. Zero formula and its branch normalization; local variable 1728/\allowbreak{}j is reciprocal to the family t.\newline Assembly with supplied source and review context; not an independent proof review.\newline \path{records/periods/citation_checks.json}; \url{https://www.math.ucla.edu/~wdduke/preprints/zeros.pdf}\\
V113 & GPT-6 Astra, ChatGPT Pro mode\newline 2026-11-16; S\newline DM20; Theorem 1.1; Corollary 1.2, preprint pp.2--3 & opened and matched. Adelically metrized height line and small-height Galois-orbit probability equidistribution for a section of nonzero generic canonical height over a number field.\newline Assembly with supplied source and review context; not an independent proof review.\newline \path{records/periods/citation_checks.json}; \url{https://arxiv.org/pdf/1701.07947}\\
V114 & GPT-6 Astra, ChatGPT Pro mode\newline 2026-11-16; S\newline Johansson2016; §§2--3 & opened and matched. Midpoint-radius representation and outward error bounds underpin the supplied Arb computation.\newline Assembly with supplied source and review context; not an independent proof review.\newline \path{records/periods/citation_checks.json}; \url{https://arxiv.org/pdf/1611.02831}\\
V115 & GPT-6 Astra, ChatGPT Pro mode\newline 2026-11-16; S\newline PythonFlint2026; arb constructor, abs\_\allowbreak{}upper and upper methods & opened and matched. Exact dyadic/\allowbreak{}rational constructor inputs, ball enclosures, outward upper endpoints.\newline Assembly with supplied source and review context; not an independent proof review.\newline \path{records/periods/citation_checks.json}; \url{https://python-flint.readthedocs.io/en/latest/arb.html}\\
V116 & GPT-6 Astra, ChatGPT Pro mode\newline 2026-11-16; S\newline DGKXv2; Abstract and Introduction, lines 134--270; whole-file searches for p2:\allowbreak{}thm:\allowbreak{}exact-period and p2:\allowbreak{}cor:\allowbreak{}bare-period & opened and matched. The exact-period theorem present in version 1 is omitted in version 2. No mathematical conclusion is inferred from this omission.\newline Assembly with supplied source and review context; not an independent proof review.\newline \path{records/periods/citation_checks.json}; \path{papers/arXiv_2610.00250v2_Du_Guo_Kusakabe_Xie.tex}\\
V117 & GPT-6 Astra, ChatGPT Pro mode\newline 2026-11-16; S\newline AC26; Supplied transcription:\allowbreak{} Definition 3.1 and Theorem 3.4(i), printed p. 14; Theorem 9.1, p. 64; Proposition 9.11, pp. 70--71; Lemma 9.12, pp. 72--74; Proposition 9.13, pp. 74--76; Remarks 1.3 and 9.21, pp. 5 and 81; Corollary 9.22 and Proposition 9.23, p. 81; Theorems 7.17 and 7.22; Lemma 10.3 and Theorem 10.5, pp. 87--89. & opened and matched. Opened and matched construction/\allowbreak{}period inputs, known structure direct images and Hodge constraints, canonical and topology credits, correction targets, and the memberwise algebraic-dimension input. No source prose reproduced. Verification concerns the supplied transcription, not a newly inspected original PDF.\newline Assembly with supplied source and review context; not an independent proof review.\newline \path{records/cohomology/citation_checks.json}; \url{https://alpo.ge/s6.pdf}\\
V118 & GPT-6 Astra, ChatGPT Pro mode\newline 2026-11-16; S\newline E26; Introduction, lines 215--277; proposition fundamental-group, lines 1215--1231; final invariant-Hodge-form remark, lines 1723--1731. & opened and matched. Opened the original supplied TeX; it credits AC26 for Dolbeault information and expressly omits those calculations. The construction comparison, standard-collar presentation and fibre-form background match the cited/\allowbreak{}credited scope.\newline Assembly with supplied source and review context; not an independent proof review.\newline \path{records/cohomology/citation_checks.json}; \url{https://arxiv.org/abs/2609.38442}\\
V119 & GPT-6 Astra, ChatGPT Pro mode\newline 2026-11-16; S\newline DGKXv1; Theorem 1.3 (p2:\allowbreak{}thm:\allowbreak{}kuranishi-aut), lines 310--327; Proposition 6.10 (p2:\allowbreak{}prop:\allowbreak{}structure-direct-images), lines 9151ff.; Theorem 6.12 (p2:\allowbreak{}thm:\allowbreak{}tangent-cohomology), lines 9445ff.; p2:\allowbreak{}prop:\allowbreak{}kuranishi-line, lines 9877ff.; p2:\allowbreak{}thm:\allowbreak{}intrinsic-fibration, lines 10018ff. & opened and matched. The exact tangent vector, smooth reduced Kuranishi line and miniversality are already present. The structure direct images and intrinsic-pencil inputs also match. The reviewed two DGKX lemma citations have been replaced by BS76, RRV71 and RR70 rather than retained as primary analytic references.\newline Assembly with supplied source and review context; not an independent proof review.\newline \path{records/cohomology/citation_checks.json}; \url{https://arxiv.org/abs/2610.00250v1}\\
V120 & GPT-6 Astra, ChatGPT Pro mode\newline 2026-11-16; S\newline VL26; Theorem 1.1 (t:\allowbreak{}main), source lines 210--244; Proposition 2.3 (logprop), lines 420--496. & opened and matched. Opened the actual source theorem and logarithmic-transformation proposition. They provide the sufficiently small nonzero construction domain, coprimality and admissible free cyclic quotients required in inv:\allowbreak{}B2.1.\newline Assembly with supplied source and review context; not an independent proof review.\newline \path{records/cohomology/citation_checks.json}; \url{https://arxiv.org/abs/2610.39102}\\
V121 & GPT-6 Astra, ChatGPT Pro mode\newline 2026-11-16; S\newline BS76; Chapter III, Theorem 4.10, pp. 132--133; Theorem 4.12(ii), pp. 134--135; web lines 5410--5508. & opened and matched. Original-book text reproduction opened. The first theorem supplies the analytic base-change criterion and comparison; the second gives cohomological flatness and local freeness from constant fibre-cohomology dimension for a flat coherent sheaf over a reduced base. Applied on the source complex-space charts, with canonical overlap comparison.\newline Assembly with supplied source and review context; not an independent proof review.\newline \path{records/cohomology/citation_checks.json}; \url{https://djvu.online/file/2hQiOnz8WuoQk}\\
V122 & GPT-6 Astra, ChatGPT Pro mode\newline 2026-11-16; S\newline RRV71; Original article introduction p. 261 and Section 2, pp. 264--266; original bibliographic identification Invent. Math. 13 (1971), 261--283, DOI 10.1007/\allowbreak{}BF01406078. & opened and matched. Full original-article text opened through a third-party reproduction. The proper analytic duality isomorphism and relative trace were matched. These support the source duality applications on smooth analytic atlases; RR70 identifies the Gorenstein dualizing line.\newline Assembly with supplied source and review context; not an independent proof review.\newline \path{records/cohomology/citation_checks.json}; \url{https://www.academia.edu/5435625/Dualit%C3%A9_relative_en_g%C3%A9om%C3%A9trie_analytique_complexe}\\
V123 & GPT-6 Astra, ChatGPT Pro mode\newline 2026-11-16; S\newline RR70; Section 3, Theorem 1, printed p. 81; Proposition 2 and complete-intersection paragraph, printed p. 91 (PDF page 16). & opened and matched. Original Numdam PDF opened. A Gorenstein analytic space has its dualizing complex represented by an invertible sheaf in the dimension shift; complete intersections are included. This is the dualizing-sheaf input used with the relative theorem.\newline Assembly with supplied source and review context; not an independent proof review.\newline \path{records/cohomology/citation_checks.json}; \url{https://www.numdam.org/item/10.1007/BF02684652.pdf}\\
V124 & GPT-6 Astra, ChatGPT Pro mode\newline 2026-11-16; S\newline Ugarte00; Proposition 3.1, pp. 175--176; Theorem 3.2, pp. 176--177; Corollary 3.3, p. 177. & opened and matched. Original PDF opened through the already retrieved primary-source reference after two fresh URL timeouts. Matched the conditional Hodge relations and the second- and third-page Frölicher constraints for a hypothetical complex S6. The paper cites these as background, not as an evaluation of this family’s h12.\newline Assembly with supplied source and review context; not an independent proof review.\newline \path{records/cohomology/citation_checks.json}; \url{https://zaguan.unizar.es/record/165697/files/texto_completo.pdf?version=1}\\
V125 & GPT-6 Astra, ChatGPT Pro mode\newline 2026-11-16; S\newline McHugh18; Version 4, 15 November 2018; Section 7.1, printed pp. 24--25, including the table on p. 25; web lines 3342--3500. & opened and matched. Original PDF opened. The formulas express complex-S6 Bott--Chern dimensions using Dolbeault dimensions and the E2 term, with Aeppli obtained by the stated duality. Exact comparison scope is the background recorded in inv:\allowbreak{}C2.2; no certification of the whole paper is asserted.\newline Assembly with supplied source and review context; not an independent proof review.\newline \path{records/cohomology/citation_checks.json}; \url{https://arxiv.org/pdf/1708.03251}\\
V126 & GPT-6 Astra, ChatGPT Pro mode\newline 2026-11-16; S\newline Angella17; Version 2, November 2017; introduction, Figures 1--2, and Section 4, printed pp. 13--16. & opened and matched. Original PDF opened. Matched the conditional spectral-sequence/\allowbreak{}double-complex diagrams and Bott--Chern relations for complex S6. Used only for the precise historical discussion and credit, with the actual table proved from the stated INV input.\newline Assembly with supplied source and review context; not an independent proof review.\newline \path{records/cohomology/citation_checks.json}; \url{https://arxiv.org/pdf/1705.10518}\\
V127 & GPT-6 Astra, ChatGPT Pro mode\newline 2026-11-16; S\newline Stelzig21; Version 3, 8 March 2021; Theorem 3 (Theorem A in introduction, printed p. 2), Proposition 6(1), p. 16, Lemma 8, pp. 17--18. & opened and matched. Original PDF opened. The bounded-double-complex decomposition into squares and zigzags is valid for vector spaces without a finite-dimensional-component hypothesis. Proposition 6(1) identifies even-zigzag multiplicities with ranks of the spectral-sequence differentials; the BC/\allowbreak{}A counting interpretation was matched.\newline Assembly with supplied source and review context; not an independent proof review.\newline \path{records/cohomology/citation_checks.json}; \url{https://arxiv.org/pdf/1812.00865}\\
V128 & GPT-6 Astra, ChatGPT Pro mode\newline 2026-11-16; S\newline AC26; §9.7, Proposition 9.23 and Corollary 9.24, transcription lines4340--4390 & opened and matched. Aut\textasciicircum{}0=C*, its additive translation coordinate and effective period Z; the fixed double curve and semifree normal weights. Full Aut statement is credited separately to DGKXv1.\newline Assembly check with supplied proofs and review context; source opened and matched, not a new independent proof verification.\newline \path{records/geometry_ends/citation_checks.json}; \path{original/AC26_transcription.tex}\\
V129 & GPT-6 Astra, ChatGPT Pro mode\newline 2026-11-16; S\newline AC26; §9.1, Lemmas9.2--9.4, Proposition9.5 and Corollary9.8; transcription lines3274--3551 & opened and matched. Exact six-coefficient NS equation; five-function independence; indefiniteness; countable exceptional locus and generic NS=Z eta. The +12 wedge coefficient is distinguished from the complex-oriented numerical self-intersection −12.\newline Assembly check with supplied proofs and review context; source opened and matched, not a new independent proof verification.\newline \path{records/geometry_ends/citation_checks.json}; \path{original/AC26_transcription.tex}\\
V130 & GPT-6 Astra, ChatGPT Pro mode\newline 2026-11-16; S\newline AC26; §4.6 Proposition4.6, source PDF p.29; transcription lines1452--1470 & opened and matched. Finite dP6 normalization, opposite boundary-pair identifications, three double rational curves, two triple points and Euler number2.\newline Assembly check with supplied proofs and review context; source opened and matched, not a new independent proof verification.\newline \path{records/geometry_ends/citation_checks.json}; \path{original/AC26_transcription.tex}\\
V131 & GPT-6 Astra, ChatGPT Pro mode\newline 2026-11-16; S\newline E26; §2 Lemma2.1 and Corollary2.2; source lines330--401 & opened and matched. Height(P)=1/\allowbreak{}6; intersection inputs; the line O(P−O)⊗t6P*O(O−P) is the degree-one base pullback and its linearization section has divisor S\_\allowbreak{}infinity.\newline Assembly check with supplied proofs and review context; source opened and matched, not a new independent proof verification.\newline \path{records/geometry_ends/citation_checks.json}; \path{papers/arXiv_2609.38442_Engel_src/S6.tex}\\
V132 & GPT-6 Astra, ChatGPT Pro mode\newline 2026-11-16; S\newline E26; §8 Conjecture8.2, lines1673--1676, and Remarks8.1--8.4 & opened and matched. Exact as-worded conjecture for the specified nonnormal pinching; Remark8.4 asks Klingler’s question for the indefinite (1,6) domain. No stated exact index/\allowbreak{}maximal-open answer.\newline Assembly check with supplied proofs and review context; source opened and matched, not a new independent proof verification.\newline \path{records/geometry_ends/citation_checks.json}; \path{papers/arXiv_2609.38442_Engel_src/S6.tex}\\
V133 & GPT-6 Astra, ChatGPT Pro mode\newline 2026-11-16; S\newline DGKXv1; Theorem p2:\allowbreak{}thm:\allowbreak{}full-automorphisms, source line11086, and proof & opened and matched. Aut(X)=Aut\_\allowbreak{}\{P1\}(X)=C*, including every component. The relevant precise theorem is the full-automorphisms theorem, rather than the exact-period theorem alone.\newline Assembly check with supplied proofs and review context; source opened and matched, not a new independent proof verification.\newline \path{records/geometry_ends/citation_checks.json}; \path{papers/arXiv_2610.00250v1_Du_Guo_Kusakabe_Xie.tex}\\
V134 & GPT-6 Astra, ChatGPT Pro mode\newline 2026-11-16; S\newline Schnell19; Theorem4.1 p.1; equation(4.3) p.2; effective-bundle quotient discussion pp.2--3; Theorem4.4 p.3 & opened and matched. Appell--Humbert correspondence, theta factor of automorphy, semipositive effective line bundles and positive integral Hermitian criterion for projectivity. Used in FAM D.1.6 and D.2.3.\newline Assembly check with supplied proofs and review context; source opened and matched, not a new independent proof verification.\newline \path{records/geometry_ends/citation_checks.json}; \url{https://www.math.stonybrook.edu/~cschnell/mat615/lectures/lecture4.pdf}\\
V135 & GPT-6 Astra, ChatGPT Pro mode\newline 2026-11-16; S\newline Buchdahl1999; Theorem11, printed p.297 & opened and matched. A compact complex surface with even first Betti number is Kähler. Applied to smooth connected surfaces with b1=2 in PUSH D.1.2.\newline Assembly check with supplied proofs and review context; source opened and matched, not a new independent proof verification.\newline \path{records/geometry_ends/citation_checks.json}; \url{https://www.numdam.org/article/AIF_1999__49_1_287_0.pdf}\\
V136 & GPT-6 Astra, ChatGPT Pro mode\newline 2026-11-16; S\newline Grauert1960D; §6 HauptsatzI, printedp.59; §7 Satz1, printedp.60 & opened and matched. Coherence of proper higher direct images and proper analytic image theorem. The proof does not use the originally stated closed exceptional locus of Satz2.\newline Assembly check with supplied proofs and review context; source opened and matched, not a new independent proof verification.\newline \path{records/geometry_ends/citation_checks.json}; \url{https://www.numdam.org/article/PMIHES_1960__5__5_0.pdf}\\
V137 & GPT-6 Astra, ChatGPT Pro mode\newline 2026-11-16; S\newline Grauert1963ErratumD; Item2, p.36 & opened and matched. Corrects the claimed closedness of the exceptional set in the 1960 §7 Satz2. PUSH D.1.2 instead proves the needed base change over a smooth curve directly.\newline Assembly check with supplied proofs and review context; source opened and matched, not a new independent proof verification.\newline \path{records/geometry_ends/citation_checks.json}; \url{https://www.numdam.org/item/PMIHES_1963__16__35_0.pdf}\\
V138 & GPT-6 Astra, ChatGPT Pro mode\newline 2026-11-16; S\newline Levy08; Lemmas3.1--3.2; §3.1, pp.588--590 & opened and matched. Analytic complexes give topological K-classes; the local complex comparison and stabilization/\allowbreak{}descent support the virtual tangent class used in PUSH A.3.1.\newline Assembly check with supplied proofs and review context; source opened and matched, not a new independent proof verification.\newline \path{records/geometry_ends/citation_checks.json}; \url{https://www.numdam.org/article/AIF_2008__58_2_571_0.pdf}\\
V139 & GPT-6 Astra, ChatGPT Pro mode\newline 2026-11-16; S\newline Loj64; §3 Theorems1--2, pp.463--464 & opened and matched. Compatible triangulation for the specified locally finite semianalytic sets in affine space and a countable real analytic manifold; applied to the compact radial inverse image and its singular/\allowbreak{}end strata.\newline Assembly check with supplied proofs and review context; source opened and matched, not a new independent proof verification.\newline \path{records/geometry_ends/citation_checks.json}; \url{https://www.numdam.org/item/ASNSP_1964_3_18_4_449_0.pdf}\\
V140 & GPT-6 Astra, ChatGPT Pro mode\newline 2026-11-16; S\newline FG10; §5.1 Theorem5.8 and Remark5.9, p.102 & opened and matched. Virtual Euler/\allowbreak{}Chern-number and smoothing result in the algebraic scheme setting. The paper cites this only as the algebraic counterpart, expressly not as a replacement for the analytic argument.\newline Assembly check with supplied proofs and review context; source opened and matched, not a new independent proof verification.\newline \path{records/geometry_ends/citation_checks.json}; \url{https://msp.org/gt/2010/14-1/gt-v14-n1-p02-s.pdf}\\
V141 & GPT-6 Astra, ChatGPT Pro mode\newline 2026-11-16; S\newline Ehresmann52; §1 Proposition and proof, pp.154--155 & opened and matched. A full-rank differentiable surjection with compact connected fibres has differentiable bundle structure; horizontal transport identifies the compact smooth fibres in PUSH D.1.2.\newline Assembly check with supplied proofs and review context; source opened and matched, not a new independent proof verification.\newline \path{records/geometry_ends/citation_checks.json}; \url{https://www.numdam.org/item/SB_1948-1951__1__153_0.pdf}\\
V142 & GPT-6 Astra, ChatGPT Pro mode\newline 2026-11-16; S\newline Meersseman2019E; §2.1 pp.883--884; §3.2 Theorem3.4 and Remark3.6 pp.898--899 & opened and matched. The smooth integrable-structure space has connected open local Kuranishi charts and a retraction to a finite-dimensional analytic space; the local product restricts to C-infinity operators. This is the mechanism cited for PUSH E.2.24; no global bounded-automorphism hypothesis is imported.\newline S-source match with supplied proof context; no independent F or D check asserted.\newline \path{records/geometry_ends/citation_checks.json}; \url{https://www.numdam.org/item/10.5802/jep.108.pdf}\\
\bottomrule\end{longtable}\endgroup
No human has checked any part of the mathematics. No formalization was supplied. No new execution of the missing mathematical programs is claimed. Gray is not a mathematical input; the public repository link was supplied but could not be fetched. Source comparisons and build checks are recorded separately from mathematical verification.
